\chapter{Fluid dynamics: inviscid}
\label{vol03:ch11}

\epigraph{Things which are not yet known, are not on that account
to be confounded with things which are not. Pressure produced in
fluids by external causes is communicated equally in every
direction; and the velocity of efflux of a fluid at any opening is
that of a heavy body falling through a height equal to the height
of the fluid above the opening.}{Daniel Bernoulli, \emph{
Hydrodynamica} (1738), Section XII, in the Carmody and Kobus
translation \cite[pp.~244--245]{hist:bernoulli1738}.}

\chapmeta{Half-life: durable. Archetypes: balance, transport. Exercise target: 39. Project track: physical build (Pitot tube on a household fan). Hazard class: low. Reader path default: standard, with sections explicitly tagged. Named cases: Tacoma Narrows aeroelastic flutter (1940, cross-referenced to Chapter 13); Pitot-icing common-mode failure (Air France 447, cross-referenced to Chapter 12).}

\noindent
The previous chapter treated the fluid at rest: pressure as a
scalar field, hydrostatics as a special case of mechanical
equilibrium. This chapter sets the fluid in motion under the
assumption of zero viscosity. The inviscid idealisation cuts away
the thin boundary-layer shear that dominates wall drag, leaving the
bulk-flow physics that governs lift, pressure recovery, jet
trajectories, and free-surface waves. The reader leaves able to
write the continuity equation for any control volume, derive the
Euler momentum equation along a streamline, apply Bernoulli's
equation with full awareness of its assumptions, and recognise the
configurations in which a Bernoulli calculation gives the right
answer to within a few percent and those in which the same
calculation is qualitatively wrong.

\section{Continuity: mass conservation in a flowing fluid}\pathtag{core}

A flowing fluid carries mass, momentum, and energy through space.
The first conservation law is mass. Pick a fixed region of space,
a \engterm{control volume} $\mathcal{V}$ with boundary
$\partial\mathcal{V}$, and ask: how does the mass inside the
control volume change with time? The answer is the integral
continuity equation:
\[
\frac{d}{dt}\int_{\mathcal{V}} \rho \,dV
+ \oint_{\partial\mathcal{V}} \rho \vec{u} \cdot \hat{n}\,dA = 0,
\]
where $\rho$ is the local fluid density, $\vec{u}$ the local fluid
velocity, and $\hat{n}$ the outward unit normal on the boundary.
The first term is the rate of change of mass stored inside; the
second is the net mass flux leaving across the boundary. Their sum
is zero because mass is neither created nor destroyed inside the
control volume.

Applying the divergence theorem to the boundary integral and
collecting both terms under a single volume integral, the integrand
must vanish pointwise (since the control volume was arbitrary):
\[
\frac{\partial \rho}{\partial t} + \nabla\cdot(\rho\vec{u}) = 0.
\]
This is the \engterm{differential continuity equation}, valid at
every point of the flow at every instant. The two forms (integral
and differential) carry the same physical content; the integral
form is the working tool for a control-volume calculation around a
pipe, a tank, or an engine; the differential form is the working
tool for a partial-differential-equation analysis of the flow
field.

For an \engterm{incompressible} fluid, $\rho$ is constant in space
and time, and the differential continuity equation reduces to
\[
\nabla\cdot\vec{u} = 0.
\]
The velocity field of an incompressible flow is solenoidal: it has
zero divergence at every point. Liquids satisfy the incompressible
condition to a very good approximation under nearly all engineering
conditions; the bulk modulus of water is about $2.2\,\text{GPa}$,
giving a density change of about $0.05\,\%$ under a $1\,\text{MPa}$
pressure rise. Gases satisfy the incompressible condition only at
flow speeds well below the local speed of sound, the criterion of
chapter 13 of this volume; at speeds above about $0.3 M$, density
variations of more than $5\,\%$ make incompressible analysis
inaccurate.

For a steady one-dimensional flow through a pipe or channel, the
integral continuity equation simplifies further. With no time
variation and with the inlet and outlet streams normal to the
control surface,
\[
\rho_1 A_1 u_1 = \rho_2 A_2 u_2,
\]
where subscripts $1$ and $2$ label the inlet and outlet
cross-sections, $A$ is the cross-sectional area, and $u$ is the
flow speed averaged over the cross-section. For an incompressible
fluid, $\rho_1 = \rho_2$ and the relation reduces to
\[
A_1 u_1 = A_2 u_2,
\]
the \engterm{conservation of volume flow rate} $Q = A u$. A pipe
that narrows from $100\,\si{\milli\meter}$ to $50\,\si{\milli\meter}$
diameter accelerates the flow by a factor of four (the area ratio
is the square of the diameter ratio). The acceleration is purely
geometric; no force has been invoked yet. The next section supplies
the force.

% Figure: control volume in a contracting duct, with mass fluxes in and out.
% Shows the integral continuity bookkeeping: rho1 A1 u1 in, rho2 A2 u2 out.
\begin{figure}[htbp]
\centering
\begin{tikzpicture}[
  duct/.style={draw=engblack, very thick},
  cv/.style={draw=engred, thick, dashed},
  flux/.style={draw=engblue, -{Stealth}, thick},
  label/.style={font=\small},
]
% Contracting duct walls (upper and lower)
\draw[duct] (-4,1.6) -- (0,1.6) -- (3,0.8) -- (6,0.8);
\draw[duct] (-4,-1.6) -- (0,-1.6) -- (3,-0.8) -- (6,-0.8);

% Control volume boundary (dashed): inlet face at x=-3, outlet face at x=5
\draw[cv] (-3,-1.6) -- (-3,1.6);
\draw[cv] (5,-0.8) -- (5,0.8);
\draw[cv] (-3,1.6) -- (0,1.6) -- (3,0.8) -- (5,0.8);
\draw[cv] (-3,-1.6) -- (0,-1.6) -- (3,-0.8) -- (5,-0.8);

% Inlet flux arrows
\foreach \y in {-1.0,0,1.0} {
  \draw[flux] (-3.9,\y) -- (-3.1,\y);
}
% Outlet flux arrows (faster, shorter region, longer arrows)
\foreach \y in {-0.5,0,0.5} {
  \draw[flux] (5.1,\y) -- (6.1,\y);
}

% Inlet labels
\node[label, engblue] at (-3.6,1.35) {$\rho_1 A_1 u_1$};
\node[label, anchor=north] at (-3,-1.7) {inlet face $A_1$};
% Outlet labels
\node[label, engblue] at (6.0,1.0) {$\rho_2 A_2 u_2$};
\node[label, anchor=north] at (5,-0.95) {outlet face $A_2$};

% Outward normals
\draw[-{Latex}, engred, thick] (-3,0) -- (-3.7,0);
\node[engred, font=\scriptsize, anchor=east] at (-3.7,0.2) {$\hat{n}$};
\draw[-{Latex}, engred, thick] (5,0) -- (5.7,0);
\node[engred, font=\scriptsize, anchor=west] at (5.7,0.2) {$\hat{n}$};

% Storage label inside CV
\node[engred, font=\small] at (1.0,0) {$\displaystyle\frac{d}{dt}\!\int_{\mathcal{V}}\!\rho\,dV$};

\end{tikzpicture}
\caption[Control volume in a contracting duct]{The integral
continuity balance drawn on a fixed control volume (red dashed
boundary) spanning a contracting duct. Mass enters across the inlet
face at rate $\rho_1 A_1 u_1$ and leaves across the outlet face at
rate $\rho_2 A_2 u_2$; the outward normal $\hat{n}$ points out of the
volume on every face, so the inlet flux $\rho\vec{u}\cdot\hat{n}$ is
negative and the outlet flux positive. For steady flow the stored
mass $\int_{\mathcal{V}}\rho\,dV$ does not change, so the two face
fluxes balance: $\rho_1 A_1 u_1 = \rho_2 A_2 u_2$. The side walls
carry no flux because the velocity there is tangent to the wall. The
acceleration through the contraction is purely geometric; no force has
been invoked.}
\label{fig:vol03ch11:control-volume}
\end{figure}


Figure~\ref{fig:vol03ch11:control-volume} draws the bookkeeping on a
contracting duct: mass enters one face, leaves the other, and (in
steady flow) nothing accumulates between. The reader who can sketch
this control volume and write the two face fluxes has the whole of
mass conservation in hand; everything later in the chapter rests on
closing this balance first.

\begin{principle}[Continuity is a bookkeeping identity]
Mass entering a control volume per unit time equals mass leaving per
unit time plus mass accumulating inside. The relation is exact for
any control volume in any flow; it requires no assumption about the
constitutive behaviour of the fluid. For an incompressible flow
through a pipe, the velocity scales inversely with the
cross-sectional area; for a compressible flow, the density carries
part of the variation, and the velocity scales with $1/(\rho A)$.
Every fluid-mechanics calculation starts by closing the mass
balance.
\end{principle}

\section{Euler's equation: momentum conservation in an inviscid fluid}\pathtag{core}

Newton's second law applied to a fluid particle of mass $\rho\,dV$
gives
\[
\rho\,dV\,\frac{D\vec{u}}{Dt} = \vec{F}_{\text{on the particle}},
\]
where $D/Dt = \partial/\partial t + \vec{u}\cdot\nabla$ is the
\engterm{material derivative}, the time rate of change following
the fluid particle (rather than at a fixed point in space). The
material derivative carries the convective acceleration
$(\vec{u}\cdot\nabla)\vec{u}$, which records the change in velocity
the particle experiences because it moves into a region where the
velocity is different.

The forces on the particle are of two kinds: \engterm{body forces}
(gravity, electromagnetic forces) that act on the volume, and
\engterm{surface forces} (pressure, viscous stress) that act on the
particle's surface. For an \engterm{inviscid fluid}, the surface
force reduces to the pressure alone: the surrounding fluid pushes
on the particle with a force per unit area equal to the local
pressure, directed inward along the surface normal. The net surface
force per unit volume on the particle is $-\nabla p$. Adding the
gravitational body force $\rho \vec{g}$ gives
\[
\rho\,\frac{D\vec{u}}{Dt} = -\nabla p + \rho\vec{g}.
\]
This is the \engterm{Euler equation}, derived by Leonhard Euler in
his 1755 \emph{Principes g\'en\'eraux du mouvement des fluides}.
The Euler equation is the inviscid limit of the Navier-Stokes
equation (chapter 12 of this volume); it is the momentum-balance
law for a fluid in which viscosity is neglected.

Expanding the material derivative,
\[
\frac{\partial \vec{u}}{\partial t}
+ (\vec{u}\cdot\nabla)\vec{u}
= -\frac{1}{\rho}\nabla p + \vec{g}.
\]
The first term on the left is the \engterm{local acceleration} (the
unsteadiness of the flow at a fixed point); the second is the
\engterm{convective acceleration} (the spatial variation the
particle samples as it moves). The convective term is nonlinear in
$\vec{u}$; this nonlinearity is the source of most of fluid
dynamics' analytical difficulty, including turbulence (chapter 12)
and shock waves (chapter 13).

For a \engterm{steady flow}, $\partial \vec{u}/\partial t = 0$, and
the Euler equation becomes a statement about the spatial structure
of the flow field. Taking the dot product of the steady Euler
equation with a unit tangent $\hat{s}$ along a streamline,
\[
(\vec{u}\cdot\nabla)\vec{u}\cdot\hat{s}
= u\,\frac{\partial u}{\partial s}
= \frac{1}{2}\frac{\partial(u^2)}{\partial s},
\]
where $s$ is arc length along the streamline. The pressure term
becomes $-(1/\rho)\,\partial p/\partial s$, and the gravity term
becomes $-g\,\partial z/\partial s$ (with $z$ measured against
gravity). Multiplying through by $\rho$ and integrating along the
streamline gives Bernoulli's equation, the subject of the next
section.

The Euler equation also has a \engterm{normal} form, taken
perpendicular to the streamline. For a streamline with local radius
of curvature $R_c$, the centripetal acceleration is $u^2/R_c$ toward
the centre of curvature, and the normal component of Euler's
equation gives
\[
\frac{\rho u^2}{R_c} = -\frac{\partial p}{\partial n}
+ \rho g_n,
\]
where $n$ is the inward normal coordinate and $g_n$ is the gravity
component along it. The relation says: a streamline curves because
the pressure gradient across it (plus any gravity component) pushes
the fluid particle inward. The pressure is higher on the outside
of a curved streamline than on the inside. This is the working
intuition behind every aerodynamic surface: a curved streamline
over the top of an airfoil corresponds to a pressure low on the
upper surface relative to the freestream, and this pressure low
times the upper surface area is what lifts the wing.

\subsection{The integral momentum theorem}

The Euler equation is the pointwise momentum balance. Integrating it
over a fixed control volume, or applying Newton's second law directly
to the mass inside, gives the \engterm{integral momentum theorem},
the working tool whenever the forces on a device matter more than the
detailed flow field:
\[
\sum \vec{F}
= \frac{d}{dt}\int_{\mathcal{V}} \rho\vec{u}\,dV
+ \oint_{\partial\mathcal{V}} \rho\vec{u}\,(\vec{u}\cdot\hat{n})\,dA.
\]
The left side collects every force on the control volume: pressure
forces on the inlet and outlet faces, the reaction force from any
solid surface cut by the boundary, gravity, and any external
restraint. The right side is the rate of change of stored momentum
plus the net momentum flux carried out across the boundary. For
steady flow with one inlet and one outlet, the relation reduces to
\[
\sum \vec{F} = \dot{m}\,(\vec{u}_2 - \vec{u}_1),
\quad \dot{m} = \rho A u,
\]
the form a working engineer uses to size the anchor bolts on a pipe
bend, the thrust of a jet, or the force a water jet exerts on a
turbine bucket. The momentum theorem makes no inviscid assumption: it
holds across a hydraulic jump, through a separated wake, and inside a
turbomachine, the configurations where Bernoulli does not. That
contrast, energy not conserved but momentum always conserved, is the
diagnostic backbone of the failure section.

% Figure: control-volume momentum balance on a 90-degree pipe bend,
% showing the reaction force on the bend.
\begin{figure}[htbp]
\centering
\begin{tikzpicture}[
  wall/.style={draw=engblack, very thick},
  cv/.style={draw=engred, thick, dashed},
  mom/.style={draw=engblue, -{Stealth}, very thick},
  force/.style={draw=englime, -{Stealth}, very thick},
  label/.style={font=\small},
]
% Pipe bend: horizontal inlet (left), turns up (top). Outer and inner walls.
% Inlet pipe from left, vertical pipe going up.
\draw[wall] (-4,0.6) -- (-0.6,0.6) -- (-0.6,4);   % outer/upper-left into vertical right wall
\draw[wall] (-4,-0.6) -- (0.6,-0.6) -- (0.6,4);   % lower then vertical left wall
% Control-volume faces
\draw[cv] (-3.4,0.6) -- (-3.4,-0.6);             % inlet face
\node[label, anchor=east, engred] at (-3.5,0) {inlet};
\draw[cv] (-0.6,3.4) -- (0.6,3.4);               % outlet face
\node[label, anchor=south, engred] at (0,3.5) {outlet};
\draw[cv] (-3.4,0.6) -- (-0.6,0.6) -- (-0.6,3.4); % cv boundary along outer
\draw[cv] (-3.4,-0.6) -- (0.6,-0.6) -- (0.6,3.4); % cv boundary along inner

% Inlet momentum flux (rightward)
\draw[mom] (-3.9,0) -- (-2.8,0);
\node[label, engblue, anchor=south] at (-3.35,0.05) {$\dot{m}u_1$};
% Outlet momentum flux (upward)
\draw[mom] (0,3.8) -- (0,4.9);
\node[label, engblue, anchor=west] at (0.05,4.4) {$\dot{m}u_2$};

% Reaction force on bend (down-left, opposing the momentum change up-right):
\draw[force] (-0.0,0.0) -- (-1.6,-1.6);
\node[label, englime, anchor=north west] at (-1.2,-1.5) {$\vec{R}$ (force on fluid)};

% Pressure-force labels at faces
\node[label, anchor=north] at (-3.4,-0.7) {$p_1 A_1$};
\node[label, anchor=west] at (0.7,3.4) {$p_2 A_2$};

\end{tikzpicture}
\caption[Momentum balance on a pipe bend]{Control-volume momentum
balance on a $90^\circ$ pipe bend. Fluid enters the dashed control
volume moving to the right carrying momentum flux $\dot{m}u_1$ and
leaves moving upward carrying $\dot{m}u_2$; the inlet and outlet
pressure forces $p_1 A_1$ and $p_2 A_2$ act on the two open faces.
The vector momentum equation
$\sum\vec{F} = \dot{m}(\vec{u}_2 - \vec{u}_1)$ requires a net force,
supplied by the bend pushing on the fluid (force $\vec{R}$). By
Newton's third law the fluid pushes back on the bend with an equal
and opposite force, which the anchor bolts must carry. The momentum
balance needs no inviscid assumption and holds even where Bernoulli
fails, which is why it is the right tool across a hydraulic jump or a
separated wake.}
\label{fig:vol03ch11:momentum-bend}
\end{figure}


\paragraph{Worked example: force on a $90^\circ$ pipe bend.} Water
($\rho = 1000\,\si{\kilogram\per\meter\cubed}$) flows at
$3\,\si{\meter\per\second}$ through a $100\,\si{\milli\meter}$
diameter bend that turns the flow from horizontal to vertical
(Figure~\ref{fig:vol03ch11:momentum-bend}). The gauge pressure at the
inlet is $200\,\text{kPa}$; for an inviscid bend of constant area the
speed is unchanged and, taking the bend short enough to neglect the
elevation change, the pressure is unchanged through the bend, so
$u_1 = u_2 = 3\,\si{\meter\per\second}$ and $p_2 = p_1 =
200\,\text{kPa}$. The area is $A = \pi(0.05)^2 = 7.85\times
10^{-3}\,\si{\meter}^2$ and the mass flow is $\dot{m} = \rho A u =
1000\times 7.85\times 10^{-3}\times 3 = 23.6\,\si{\kilogram\per
\second}$. Take $x$ along the inlet flow and $y$ along the outlet
flow. The $x$-momentum balance reads $p_1 A + R_x = \dot{m}(0 - u_1)$,
giving $R_x = -p_1 A - \dot{m}u_1 = -(200\,000)(7.85\times 10^{-3}) -
(23.6)(3) = -1571 - 71 = -1642\,\si{\newton}$. The $y$-balance reads
$-p_2 A + R_y = \dot{m}(u_2 - 0)$, giving $R_y = p_2 A + \dot{m}u_2 =
1571 + 71 = 1642\,\si{\newton}$. The force the bolts must carry is the
reaction $-\vec{R}$, of magnitude $\sqrt{R_x^2 + R_y^2} =
2320\,\si{\newton}$ directed back along the bisector of the bend. The
pressure term dominates the momentum term by a factor of about
twenty; a designer who sizes the restraint from the momentum change
alone underestimates the load by that factor.

\section{Bernoulli's equation: the energy identity along a streamline}\pathtag{core}

Integrating the steady Euler equation along a streamline, with
$\rho$ constant (incompressible flow),
\[
\int_1^2 \left(u\,\frac{\partial u}{\partial s}
+ \frac{1}{\rho}\frac{\partial p}{\partial s}
+ g\frac{\partial z}{\partial s}\right) ds = 0,
\]
which yields the \engterm{Bernoulli equation}
\[
\frac{1}{2}u^2 + \frac{p}{\rho} + gz = \text{constant along a
streamline}.
\]
In Bernoulli's original presentation in the \emph{Hydrodynamica} of
1738, the relation was written in terms of velocity heads and
pressure heads
\cite[Section XII]{hist:bernoulli1738}; the form above is the
modern energy-density version.

The relation says: along a streamline, the sum of kinetic energy
per unit mass, pressure-work per unit mass, and gravitational
potential energy per unit mass is conserved. The Bernoulli sum is
the mechanical energy density of the fluid particle as it travels
along its streamline. Energy can convert between kinetic and
pressure forms (the fluid speeds up as the pressure drops, and
vice versa), but the total along any one streamline is the same at
every point.

Multiplying through by $\rho$ gives the \engterm{pressure form} of
Bernoulli's equation:
\[
\frac{1}{2}\rho u^2 + p + \rho g z = p_0 = \text{const.\ along a
streamline}.
\]
The constant $p_0$ is the \engterm{total pressure} or \engterm{
stagnation pressure}: the pressure that a fluid particle would have
if brought to rest adiabatically and reversibly at the same elevation.
The term $\tfrac{1}{2}\rho u^2$ is the \engterm{dynamic pressure}
$q$, and $p$ is the \engterm{static pressure}.

Dividing through by $\rho g$ gives the \engterm{head form}, in
which each term has units of length:
\[
\frac{u^2}{2g} + \frac{p}{\rho g} + z = H = \text{const.\ along a
streamline},
\]
with $H$ the \engterm{total head}. The head form is the working
language of hydraulic engineering; canal slopes, pump duty curves,
and dam spillway calculations are routinely written in heads of
water.

The assumptions on which Bernoulli's equation rests are five, and
the working engineer carries the list in active memory:

\begin{itemize}[leftmargin=*]
  \item \engterm{Steady flow}: $\partial\vec{u}/\partial t = 0$
    at every point. Pulsating flows, starting transients, and
    waves violate this assumption.
  \item \engterm{Incompressible}: $\rho$ constant along the
    streamline. Gases at Mach $> 0.3$ violate this; the
    corresponding extension is the compressible Bernoulli equation
    of chapter 13 of this volume.
  \item \engterm{Inviscid}: zero viscous stress. Real fluids in
    real geometries develop boundary layers (chapter 12 of this
    volume) where viscous effects dominate; Bernoulli does not
    apply inside the boundary layer.
  \item \engterm{Along a streamline}: the constant differs between
    different streamlines unless the flow is also irrotational
    (next section).
  \item \engterm{No heat addition, no shaft work, no separation}:
    the streamline must run from inlet to outlet through a
    well-defined flow without recirculation or energy exchange.
\end{itemize}

When all five assumptions are met, Bernoulli's equation is exact.
When one or more are violated, the equation is either approximate
(the most common case) or qualitatively wrong (the case studied in
the failure section of this chapter).

% Figure: the three Bernoulli heads (velocity, pressure, elevation) along a
% streamline through a contraction, showing total head constant.
\begin{figure}[htbp]
\centering
\begin{tikzpicture}
\begin{axis}[
  width=12cm, height=7cm,
  xlabel={position along streamline $s$ (arb.\ units)},
  ylabel={head (m of water)},
  xmin=0, xmax=10, ymin=0, ymax=11,
  axis lines=left,
  legend pos=south west,
  legend cell align=left,
  grid=major, grid style={enggray!20},
  tick label style={font=\scriptsize},
  label style={font=\small},
  legend style={font=\scriptsize},
]
% Velocity head u^2/(2g): peaks at the throat near s=5
\addplot[engblue, very thick, samples=200, domain=0:10]
  {1 + 3.2*exp(-((x-5)^2)/2.2)};
\addlegendentry{velocity head $u^2/2g$}
% Pressure head p/(rho g): drops where flow speeds up
\addplot[engred, very thick, samples=200, domain=0:10]
  {7 - 3.2*exp(-((x-5)^2)/2.2)};
\addlegendentry{pressure head $p/\rho g$}
% Total head (constant): sum of the two above (elevation taken as datum)
\addplot[engblack, dashed, very thick, samples=2, domain=0:10] {8};
\addlegendentry{total head $H = u^2/2g + p/\rho g$}
\node[engblack, font=\scriptsize, anchor=south] at (axis cs:5,8.1) {$H$ constant};
\node[engblue, font=\scriptsize, anchor=south] at (axis cs:5,4.3) {throat};
\end{axis}
\end{tikzpicture}
\caption[Bernoulli heads along a streamline]{The velocity and
pressure terms of the head form of Bernoulli's equation along a
streamline that passes through a throat near $s=5$ (elevation taken
as the datum, $z=0$). As the flow accelerates into the throat the
velocity head $u^2/2g$ (blue) rises and the pressure head $p/\rho g$
(red) falls by the same amount, so their sum, the total head $H$
(dashed), stays constant. Energy converts between pressure and
kinetic forms without loss; the inviscid, steady, incompressible
assumptions are what keep the dashed line flat. A real duct would
show $H$ sloping gently downward as friction removes mechanical
energy. Computed from a Gaussian throat model.}
\label{fig:vol03ch11:bernoulli-streamline}
\end{figure}


Figure~\ref{fig:vol03ch11:bernoulli-streamline} makes the trade
visible: along a streamline through a throat, the velocity head and
the pressure head exchange magnitude while their sum, the total head,
holds flat. The flat total head is the geometric content of
Bernoulli's equation. A real duct with friction would show the total
head sloping gently downward; reading the slope of the total-head
line is how a hydraulic engineer separates the inviscid prediction
from the viscous loss.

\begin{principle}[Bernoulli is conservation of mechanical energy along a streamline]
For a steady, incompressible, inviscid flow, the quantity
$\tfrac{1}{2}\rho u^2 + p + \rho g z$ is conserved along any
streamline. The pressure drops where the fluid speeds up and rises
where the fluid slows; the elevation enters as a gravitational
potential. The relation is the working tool for jet velocity,
Pitot-tube anemometry, Venturi flow metering, pump performance, and
free-surface flow analysis.
\end{principle}

\subsection{Three standard applications}

\engterm{Torricelli's law}. A tank of liquid with a small hole at
depth $h$ below the free surface drains through the hole at speed
\[
u = \sqrt{2 g h},
\]
the same speed a body would acquire falling freely through height
$h$. The result follows from Bernoulli along a streamline from the
free surface (where $u_1 \approx 0$ by the assumption that the tank
is much larger than the hole) to the jet (where $p_2 \approx p_1$
because both connect to the atmosphere), with $z_1 - z_2 = h$. The
result was first stated by Evangelista Torricelli in 1644 and is
sometimes used as the starting axiom of fluid dynamics; Bernoulli's
derivation places it within a wider conservation principle. The
efflux speed $\sqrt{2gh}$ is the instantaneous value; when the tank
drains, $h$ falls and the problem becomes unsteady. Volume
conservation, $A_s\,dh/dt = -C_d A_o\sqrt{2gh}$, integrates to a drain
time $t = (A_s/C_d A_o)\sqrt{2h_0/g}$. The script
\path{code/tank_drain.py} integrates this ODE and checks it against
the closed form, writing the depth history to
\path{data/tank_drain_history.csv}; for a $25\,\si{\meter}\times
10\,\si{\meter}$ pool $2\,\si{\meter}$ deep draining through a
$100\,\si{\milli\meter}$ outlet, the two agree at about nine hours.

\engterm{Pitot tube}. A Pitot tube has a forward-facing opening
that registers the stagnation pressure $p_0$ and a static-pressure
port that registers $p$. The difference is the dynamic pressure
$q = p_0 - p = \tfrac{1}{2}\rho u^2$, from which
\[
u = \sqrt{\frac{2(p_0 - p)}{\rho}}.
\]
Pitot-static probes on aircraft, on wind-tunnel test sections, and
on ship hulls all rest on this relation. The probe's calibration
adjustment is typically a small geometric correction; the working
physics is Bernoulli.

\engterm{Venturi meter}. A pipe with a contraction from area $A_1$
to area $A_2$ and back accelerates the flow through the throat
and decelerates it back to the original speed. Continuity gives
$u_2 = u_1 (A_1/A_2)$; Bernoulli along the centreline gives
$p_1 - p_2 = \tfrac{1}{2}\rho(u_2^2 - u_1^2)$. Combining,
\[
u_1 = \sqrt{\frac{2(p_1 - p_2)}{\rho [(A_1/A_2)^2 - 1]}}.
\]
The flow rate $Q = A_1 u_1$ follows. The Venturi meter measures
flow rate by measuring a pressure differential; the relation is
calibrated with a discharge coefficient $C_d \approx 0.97$ that
accounts for small viscous losses (the topic of chapter 12 of this
volume). Orifice plates use the same physics with a discharge
coefficient typically in the range $0.60$ to $0.65$.

% Figure: Venturi meter (left) and Pitot-static probe (right), both reading
% a pressure difference and converting it to a flow speed.
\begin{figure}[htbp]
\centering
\begin{tikzpicture}[
  wall/.style={draw=engblack, very thick},
  stream/.style={draw=engblue, -{Stealth}, thick},
  tap/.style={draw=engred, thick},
  label/.style={font=\scriptsize},
]
% ============ Venturi meter (left) ============
\begin{scope}[xshift=0cm]
% Upper and lower walls of a Venturi: wide, throat, wide
\draw[wall] (-3,1.2) -- (-1,1.2) -- (-0.2,0.45) -- (0.2,0.45) -- (1,1.2) -- (3,1.2);
\draw[wall] (-3,-1.2) -- (-1,-1.2) -- (-0.2,-0.45) -- (0.2,-0.45) -- (1,-1.2) -- (3,-1.2);
% Streamline
\draw[stream] (-2.8,0) -- (-1.2,0);
\draw[stream] (0.5,0) -- (2.8,0);
% Pressure taps
\draw[tap] (-2,1.2) -- (-2,2.0);
\draw[tap] (0,0.45) -- (0,2.0);
\node[label, engred, anchor=south] at (-2,2.0) {$p_1$};
\node[label, engred, anchor=south] at (0,2.0) {$p_2$};
% Manometer connection
\draw[tap] (-2,2.0) -- (0,2.0);
\node[label] at (-1,2.2) {$p_1 - p_2$};
% Section labels
\node[label, anchor=north] at (-2,-1.3) {$A_1$};
\node[label, anchor=north] at (0,-0.55) {$A_2$ (throat)};
\node[font=\small, anchor=north] at (0,-2.2) {Venturi meter};
\node[label] at (0,-1.7) {$u_1=\sqrt{\dfrac{2(p_1-p_2)}{\rho[(A_1/A_2)^2-1]}}$};
\end{scope}

% ============ Pitot-static probe (right) ============
\begin{scope}[xshift=8.5cm]
% Freestream arrows
\foreach \y in {-0.7,0,0.7} {
  \draw[stream] (-3,\y) -- (-2.0,\y);
}
\node[label, engblue, anchor=south] at (-2.5,0.75) {$U_\infty$};
% Probe body: outer tube
\draw[wall] (-1.6,0.28) -- (1.6,0.28) -- (1.6,0.9);
\draw[wall] (-1.6,-0.28) -- (1.6,-0.28) -- (1.6,-0.9);
% Rounded nose (stagnation point)
\draw[wall] (-1.6,0.28) .. controls (-2.0,0.15) and (-2.0,-0.15) .. (-1.6,-0.28);
\node[label, engred, anchor=east] at (-2.0,0) {stag.};
\filldraw[engred] (-1.95,0) circle (1.3pt);
% Inner stagnation tube to total-pressure port
\draw[tap] (1.6,0.55) -- (2.4,0.55);
\node[label, engred, anchor=west] at (2.4,0.55) {$p_0$ total};
% Static port (side hole)
\filldraw[engblack] (-0.4,0.28) circle (1.0pt);
\draw[tap] (-0.4,0.28) .. controls (0.6,1.3) and (1.4,1.3) .. (1.6,0.9);
\draw[tap] (1.6,0.9) -- (2.4,0.9);
\node[label, engred, anchor=west] at (2.4,0.9) {$p$ static};
\node[font=\small, anchor=north] at (-0.4,-1.6) {Pitot-static probe};
\node[label] at (-0.4,-1.1) {$U_\infty=\sqrt{\dfrac{2(p_0-p)}{\rho}}$};
\end{scope}

\end{tikzpicture}
\caption[Venturi meter and Pitot-static probe]{Two Bernoulli
instruments. Left: a Venturi meter contracts the pipe from area
$A_1$ to a throat $A_2$, accelerating the flow and dropping the
pressure; continuity sets $u_2 = u_1(A_1/A_2)$ and Bernoulli converts
the measured differential $p_1 - p_2$ into the inlet speed, hence the
flow rate $Q = A_1 u_1$. Right: a Pitot-static probe registers the
stagnation pressure $p_0$ at a forward-facing port (where the flow is
brought to rest) and the static pressure $p$ at a side port; their
difference is the dynamic pressure $\tfrac{1}{2}\rho U_\infty^2$.
Both instruments measure a velocity by measuring a pressure
difference, and both depend on the static port reading the
undisturbed static pressure.}
\label{fig:vol03ch11:venturi-pitot}
\end{figure}


Figure~\ref{fig:vol03ch11:venturi-pitot} sets the Venturi meter and
the Pitot-static probe side by side; both turn a velocity measurement
into a pressure measurement, and both fail in the same way, through a
static port that does not read the true static pressure. The script
\path{code/venturi_flow.py} computes the flow rate from a measured
differential and sweeps the differential pressure to produce the
calibration table \path{data/venturi_calibration.csv}; for the
$100/40\,\si{\milli\meter}$ meter of the chapter exercise a
$30\,\text{kPa}$ differential corresponds to a throat speed of
$7.85\,\si{\meter\per\second}$ and a flow rate of $9.57\,\text{L/s}$.

\section{Streamlines, streaklines, pathlines}\pathtag{standard}

A flow field is a vector field $\vec{u}(\vec{x}, t)$. Three
geometric constructs describe the field's structure:

\engterm{Streamline}. A curve that is everywhere tangent to the
velocity field at a given instant. The differential equation of a
streamline at time $t_0$ is
\[
\frac{dx}{u_x(\vec{x}, t_0)} = \frac{dy}{u_y(\vec{x}, t_0)} =
\frac{dz}{u_z(\vec{x}, t_0)}.
\]
A streamline is a snapshot of the flow direction; streamlines
cannot cross (except at stagnation points, where the velocity is
zero).

\engterm{Pathline}. The trajectory of a specific fluid particle as
it moves through the flow over time. The differential equation is
\[
\frac{d\vec{x}_p}{dt} = \vec{u}(\vec{x}_p, t), \quad \vec{x}_p(0) =
\vec{x}_0,
\]
giving the particle's position $\vec{x}_p(t)$ as a function of time
from its initial position $\vec{x}_0$.

\engterm{Streakline}. The locus, at one instant, of all fluid
particles that have passed through a given point in the flow. A
plume of dye injected at a fixed location and photographed at one
instant traces a streakline.

For a \engterm{steady flow}, all three coincide: streamlines,
pathlines, and streaklines are the same curves in space, because
the flow field at any later time is the same as the flow field
at the initial time. For an unsteady flow they diverge, and the
distinction matters for interpreting flow visualisation.

% Figure: streamlines, pathline, streakline in an unsteady flow, showing they
% differ when the flow is time-dependent.
\begin{figure}[htbp]
\centering
\begin{tikzpicture}[
  axis/.style={draw=engblack, -{Stealth}, thick},
  sline/.style={draw=enggray, thick},
  path/.style={draw=engblue, very thick},
  streak/.style={draw=engred, very thick, dashed},
  label/.style={font=\scriptsize},
]
\draw[axis] (-0.4,0) -- (9,0) node[anchor=west, font=\small] {$x$};
\draw[axis] (0,-2.4) -- (0,2.4) node[anchor=south, font=\small] {$y$};

% Instantaneous streamlines (snapshot, nearly horizontal at this instant)
\foreach \y in {-2,-1,1,2} {
  \draw[sline] (0.3,\y) -- (8.5,\y);
}
\node[label, enggray, anchor=west] at (8.6,2) {streamlines (snapshot)};

% Pathline of one particle in oscillating cross-flow:
% x grows steadily, y oscillates. (single particle from origin)
\draw[path] plot[domain=0:8, samples=120]
  ({\x}, {0.9*sin(\x*60)});
\node[label, engblue, anchor=south west] at (5.0,0.95) {pathline (one particle)};
\filldraw[engblue] (0,0) circle (1.3pt);

% Streakline: locus of particles released from origin, at one instant.
% Phase-shifted relative to the pathline.
\draw[streak] plot[domain=0:8, samples=120]
  ({\x}, {0.9*sin(\x*60 + 120)});
\node[label, engred, anchor=north west] at (5.0,-1.0) {streakline (dye trace)};

% Injection point
\node[label, anchor=north east] at (0,-0.1) {release point};

\end{tikzpicture}
\caption[Streamlines, pathline, and streakline in unsteady
flow]{The three flow-visualisation curves for a time-dependent flow
$u_x = U_0$, $u_y = V_0\cos(\omega t)$ with steady streamwise drift
and an oscillating cross-flow. The instantaneous streamlines (grey)
are tangent to the velocity field at one frozen instant. The pathline
(blue) is the trajectory of a single particle released at the origin.
The streakline (red dashed) is the locus, at one instant, of all
particles that have passed through the release point; it is the curve
a dye filament traces. The three curves differ because the flow is
unsteady. For a steady flow they collapse onto one another, and that
coincidence is what makes steady-flow visualisation easy to read.}
\label{fig:vol03ch11:streamline-types}
\end{figure}


Figure~\ref{fig:vol03ch11:streamline-types} shows the three curves
pulling apart in a flow with steady drift and an oscillating
cross-component. A wind-tunnel smoke trace is a streakline, a tracked
particle is a pathline, and a computed instantaneous tangent field is
a set of streamlines; reading a smoke photograph as if it showed
streamlines is a standard error in unsteady flow, and it is one the
reader should be able to catch.

\engterm{Stream function (two-dimensional)}. For a
two-dimensional incompressible flow, the continuity equation
$\partial u_x/\partial x + \partial u_y/\partial y = 0$ is
satisfied identically by introducing a scalar \engterm{stream
function} $\psi(x, y)$ such that
\[
u_x = \frac{\partial \psi}{\partial y}, \quad
u_y = -\frac{\partial \psi}{\partial x}.
\]
Lines of constant $\psi$ are streamlines: along a curve where
$\psi$ does not change, the velocity component normal to the curve
is zero, so the curve is everywhere tangent to the flow. The flow
rate between two streamlines is the difference of their stream
function values, $Q = \psi_2 - \psi_1$. The stream function is the
working tool for two-dimensional incompressible flow analysis; it
turns a vector-field problem into a scalar-field problem with one
fewer dependent variable.

For an \engterm{axisymmetric flow} (cylindrical coordinates $r,
\theta, z$ with no $\theta$-dependence and $u_\theta = 0$), the
Stokes stream function $\psi(r, z)$ satisfies
\[
u_z = \frac{1}{r}\frac{\partial \psi}{\partial r}, \quad
u_r = -\frac{1}{r}\frac{\partial \psi}{\partial z},
\]
and again streamlines are lines of constant $\psi$.

\section{Vorticity and circulation}\pathtag{core}

The local rotation rate of a fluid particle is the \engterm{
vorticity}
\[
\vec{\omega} = \nabla\times\vec{u}.
\]
Vorticity is twice the angular velocity of the fluid particle
about its centre. A flow with zero vorticity at every point is
\engterm{irrotational}; a flow with non-zero vorticity is
\engterm{rotational}.

For an irrotational flow, the velocity field is the gradient of a
scalar potential:
\[
\vec{u} = \nabla\phi,
\]
where $\phi$ is the \engterm{velocity potential}. The combination
of incompressibility ($\nabla\cdot\vec{u} = 0$) and
irrotationality ($\vec{u} = \nabla\phi$) gives
\[
\nabla^2\phi = 0,
\]
Laplace's equation. Solutions of Laplace's equation are
\engterm{harmonic functions}, the same class that governs
electrostatic potentials, steady heat conduction in a homogeneous
solid, and gravitational potentials in source-free regions. The
mathematical correspondence is the reason that techniques developed
for one field translate directly to the others: the method of
images for electrostatics carries over to the method of images for
potential flow; complex-variable methods that map heat-conduction
problems map fluid-flow problems with the same conformal-mapping
machinery.

The \engterm{circulation} around a closed curve $C$ is the line
integral
\[
\Gamma = \oint_C \vec{u}\cdot d\vec{\ell}.
\]
By Stokes's theorem, $\Gamma$ equals the flux of vorticity through
any surface bounded by $C$:
\[
\Gamma = \int_S \vec{\omega}\cdot \hat{n}\,dA.
\]
For an irrotational flow in a simply connected region, the
circulation around any closed curve is zero. For an irrotational
flow in a multiply connected region (one with a hole through it,
for example a curve enclosing an airfoil), the circulation around
a curve enclosing the hole can be non-zero; the value is the same
for any curve that encloses the same set of holes, by the no-vorticity
condition in the rest of the region. The circulation around the
airfoil is what produces its lift.

\engterm{Kelvin's circulation theorem} (Lord Kelvin, 1869): for an
inviscid, incompressible (or barotropic) flow under conservative
body forces, the circulation around a material curve (a curve made
up of the same fluid particles at every instant) is conserved as
the curve moves with the flow:
\[
\frac{D\Gamma}{Dt} = 0.
\]
The theorem has a striking consequence: if a flow starts from rest
(zero circulation everywhere), it remains everywhere irrotational
unless viscosity or compressibility or non-conservative forces
introduce vorticity. In particular, an airfoil placed in a uniform
freestream and inviscid flow cannot develop circulation. The lift
of a real airfoil exists because viscosity, confined to a thin
boundary layer at the trailing edge, sheds a \engterm{starting
vortex} of opposite circulation as the wing accelerates from rest,
leaving the wing with a bound circulation equal and opposite to the
shed circulation, by Kelvin's theorem applied to the original
material curve that surrounded both the wing and the eventual shed
vortex.

\engterm{Helmholtz's vortex theorems} (Hermann von Helmholtz,
1858):

\begin{itemize}[leftmargin=*]
  \item Vortex lines (lines tangent to $\vec{\omega}$ at every
    point) move with the fluid in an inviscid barotropic flow.
  \item A vortex tube has constant strength along its length;
    vortex tubes cannot end inside a fluid, they must terminate
    on a boundary or close on themselves.
  \item The strength of a vortex tube (the circulation around a
    curve enclosing the tube once) is conserved in time.
\end{itemize}

The three theorems together govern the topology of vorticity in
inviscid flow. Smoke rings, tornado vortices, and the trailing
vortices of an aircraft wing all obey them; the wing vortices
extend downstream of the aircraft as a vortex pair, eventually
descending and merging or dissipating into atmospheric turbulence,
because the trailing vortex tubes must terminate on themselves
(no fluid-mechanical reason allows them to simply end).

\begin{mastery}
\textbf{The vorticity transport equation.} Taking the curl of the
Euler equation removes the pressure term (the curl of a gradient
vanishes) and gives the evolution of vorticity directly. For an
incompressible inviscid flow under conservative body forces,
\[
\frac{D\vec{\omega}}{Dt}
= (\vec{\omega}\cdot\nabla)\vec{u}.
\]
The left side is the rate of change of vorticity following a fluid
particle; the right side is the \engterm{vortex-stretching} term. It
has no analogue in two dimensions: when $\vec{\omega}$ and $\nabla$
lie in the same plane and $\vec{u}$ has no out-of-plane component,
$(\vec{\omega}\cdot\nabla)\vec{u} = 0$, so two-dimensional inviscid
vorticity is simply carried along unchanged, $D\omega/Dt = 0$. In
three dimensions the stretching term lets a vortex tube intensify as
it is drawn out and thinned, conserving its circulation (Helmholtz)
while raising its local vorticity. This is the mechanism by which a
bath drain spins faster as it narrows, by which a tornado concentrates
weak background rotation into a violent core, and by which turbulent
energy cascades from large eddies to small. Kelvin's circulation
theorem and the vorticity transport equation are the same physics
seen integrally and differentially: circulation around a material
loop is conserved, and the vorticity threading the loop rearranges
itself by stretching and tilting as the loop deforms. Viscosity,
absent from this inviscid statement, is what diffuses vorticity
across streamlines and ultimately dissipates it; its omission here is
exactly why an inviscid flow started from rest can never generate the
bound circulation a real wing carries without the boundary-layer
shedding of a starting vortex.
\end{mastery}

\section{Potential flow: superposition of elementary solutions}\pathtag{standard}

The linearity of Laplace's equation allows complex flows to be
constructed by superposing elementary solutions. Each elementary
solution corresponds to a basic flow type; their superposition
gives the flow around bodies of practical interest.

% Figure: superposition of elementary potential flows. Four panels:
% uniform, source, vortex, and uniform+doublet (cylinder).
\begin{figure}[htbp]
\centering
\begin{tikzpicture}[
  sline/.style={draw=engblue, thick},
  body/.style={draw=engblack, fill=enggray!25, thick},
  pt/.style={draw=engred, fill=engred, circle, inner sep=1.3pt},
  label/.style={font=\scriptsize},
  arr/.style={draw=engblue, -{Stealth}, thick},
]
% ===== Panel 1: uniform flow =====
\begin{scope}[xshift=0cm,yshift=0cm]
\draw[draw=enggray, thin] (-1.5,-1.5) rectangle (1.5,1.5);
\foreach \y in {-1.1,-0.55,0,0.55,1.1} {
  \draw[arr] (-1.3,\y) -- (1.3,\y);
}
\node[label, anchor=north] at (0,-1.6) {uniform $U_\infty$};
\end{scope}

% ===== Panel 2: source =====
\begin{scope}[xshift=4cm,yshift=0cm]
\draw[draw=enggray, thin] (-1.5,-1.5) rectangle (1.5,1.5);
\foreach \a in {0,30,...,330} {
  \draw[arr] ({0.25*cos(\a)},{0.25*sin(\a)}) -- ({1.25*cos(\a)},{1.25*sin(\a)});
}
\node[pt] at (0,0) {};
\node[label, anchor=north] at (0,-1.6) {source $m$};
\end{scope}

% ===== Panel 3: point vortex =====
\begin{scope}[xshift=8cm,yshift=0cm]
\draw[draw=enggray, thin] (-1.5,-1.5) rectangle (1.5,1.5);
\foreach \r in {0.45,0.85,1.25} {
  \draw[sline] (0,0) circle (\r);
}
\draw[arr] (1.25,0) -- (1.25,0.4);
\draw[arr] (0.85,0) -- (0.85,0.3);
\node[pt] at (0,0) {};
\node[label, anchor=north] at (0,-1.6) {vortex $\Gamma$};
\end{scope}

% ===== Panel 4: uniform + doublet = cylinder =====
\begin{scope}[xshift=12cm,yshift=0cm]
\draw[draw=enggray, thin] (-1.5,-1.5) rectangle (1.5,1.5);
% streamlines bending around a circle of radius R=0.5: smooth bulged curves.
% Each curve has baseline height \yo, bulging outward by a Gaussian near x=0.
\foreach \yo in {0.55,0.85,1.2} {
  \pgfmathsetmacro{\bulge}{0.35/(\yo)}
  \draw[sline] plot[domain=-1.5:1.5, samples=80, variable=\x]
    ({\x},{ \yo + \bulge*exp(-(\x*\x)/0.45) });
  \draw[sline] plot[domain=-1.5:1.5, samples=80, variable=\x]
    ({\x},{ -\yo - \bulge*exp(-(\x*\x)/0.45) });
}
\draw[body] (0,0) circle (0.5);
\node[label, anchor=north] at (0,-1.6) {uniform $+$ doublet};
\end{scope}

\end{tikzpicture}
\caption[Elementary potential flows and their superposition]{Building
blocks of two-dimensional potential flow. Each is a solution of
Laplace's equation $\nabla^2\phi=0$, so any sum of them is also a
solution. From left: a uniform stream $U_\infty$; a source of
strength $m$ with purely radial outflow $u_r = m/2\pi r$; a point
vortex of circulation $\Gamma$ with purely tangential flow
$u_\theta = \Gamma/2\pi r$; and the superposition of a uniform stream
with a doublet, which produces the closed streamline of a circular
cylinder. Adding a vortex to the last panel breaks the up-down
symmetry and yields the Kutta-Joukowski lift. The streamlines in the
rightmost panel are sketched from $\psi = U_\infty y(1-R^2/r^2)$.}
\label{fig:vol03ch11:potential-superposition}
\end{figure}


Figure~\ref{fig:vol03ch11:potential-superposition} collects the four
building blocks the rest of this section uses. The cylinder of the
last panel is a uniform stream plus a doublet; adding a vortex to it
is the entire content of the lifting-cylinder result that follows.

\engterm{Uniform flow}. A constant velocity $U_\infty$ in the
$x$-direction has potential and stream function
\[
\phi = U_\infty x, \quad \psi = U_\infty y.
\]
The streamlines are horizontal lines; the flow has no rotation, no
sources, no sinks.

\engterm{Source / sink}. A two-dimensional source of strength $m$
at the origin has
\[
\phi = \frac{m}{2\pi}\ln r, \quad \psi = \frac{m}{2\pi}\theta,
\]
in cylindrical coordinates $(r, \theta)$. The radial velocity is
$u_r = m/(2\pi r)$, decaying with distance from the source. A
sink is a negative source. Sources and sinks are mathematical
idealisations; real fluid injectors and drains have finite size
and a complicated near-field, but the far-field is well-described
by a source/sink solution.

\engterm{Doublet}. The limit of a source and sink approaching each
other with their strengths increasing such that the product
$\kappa = m d$ remains finite. The potential is
\[
\phi = -\frac{\kappa}{2\pi}\frac{\cos\theta}{r}, \quad
\psi = \frac{\kappa}{2\pi}\frac{\sin\theta}{r}.
\]
A doublet is the analogue of an electric dipole in electrostatics.

\engterm{Point vortex}. A two-dimensional point vortex of
circulation $\Gamma$ at the origin has
\[
\phi = \frac{\Gamma}{2\pi}\theta, \quad
\psi = -\frac{\Gamma}{2\pi}\ln r,
\]
giving a purely tangential velocity $u_\theta = \Gamma/(2\pi r)$.
The flow circulates around the vortex with a speed that decays
with distance; the vorticity is concentrated at the origin (a
delta function) and zero everywhere else. The point vortex is the
two-dimensional analogue of a smoke ring's core; real vortex cores
have finite size and a more complicated radial structure.

\subsection{The cylinder in uniform flow}

Superposing a uniform flow $U_\infty$ with a doublet of strength
$\kappa = 2\pi U_\infty R^2$ gives a flow around a circular
cylinder of radius $R$. The stream function is
\[
\psi = U_\infty r \sin\theta\left(1 - \frac{R^2}{r^2}\right),
\]
which vanishes on $r = R$ (the cylinder surface is a streamline)
and approaches $\psi = U_\infty y$ at large $r$ (uniform flow at
infinity). The velocity field is
\[
u_r = U_\infty \cos\theta\left(1 - \frac{R^2}{r^2}\right), \quad
u_\theta = -U_\infty \sin\theta\left(1 + \frac{R^2}{r^2}\right).
\]
On the cylinder surface ($r = R$), $u_r = 0$ and
$u_\theta = -2 U_\infty \sin\theta$. The flow speed on the surface
is twice the freestream at the top and bottom ($\theta = \pm\pi/2$)
and zero at the front and back stagnation points ($\theta = 0,
\pi$).

By Bernoulli, the pressure coefficient
$C_p = (p - p_\infty)/(\tfrac{1}{2}\rho U_\infty^2)$ on the
surface is
\[
C_p = 1 - 4\sin^2\theta.
\]
The pressure is maximum at the stagnation points ($C_p = 1$) and
minimum at the top and bottom ($C_p = -3$). Integrating the
pressure around the cylinder gives zero net force, in either the
streamwise or transverse direction: \engterm{d'Alembert's paradox}.
A cylinder in inviscid steady flow experiences no drag and no
lift, a result Jean le Rond d'Alembert noted in 1752 and that
foreshadowed the importance of viscosity (chapter 12 of this
volume) and unsteady wake structure (chapters 12 and 13) in real
fluid drag.

% Figure: pressure coefficient Cp = 1 - 4 sin^2(theta) around a cylinder,
% inviscid, with the symmetric distribution that yields d'Alembert's paradox.
\begin{figure}[htbp]
\centering
\begin{tikzpicture}
\begin{axis}[
  width=12cm, height=7cm,
  xlabel={angle from front stagnation point $\theta$ (degrees)},
  ylabel={pressure coefficient $C_p$},
  xmin=0, xmax=180, ymin=-3.4, ymax=1.4,
  axis lines=left,
  xtick={0,45,90,135,180},
  grid=major, grid style={enggray!20},
  tick label style={font=\scriptsize},
  label style={font=\small},
  legend style={font=\scriptsize},
  legend pos=south east,
]
% Cp = 1 - 4 sin^2(theta)
\addplot[engblue, very thick, samples=200, domain=0:180]
  {1 - 4*(sin(x))^2};
\addlegendentry{$C_p = 1 - 4\sin^2\theta$}
% Zero line emphasis
\draw[engblack, thin] (axis cs:0,0) -- (axis cs:180,0);
% Mark stagnation and minimum
\node[engred, font=\scriptsize, anchor=south west] at (axis cs:2,1) {front stagnation $C_p=1$};
\filldraw[engred] (axis cs:0,1) circle (1.6pt);
\node[engred, font=\scriptsize, anchor=north] at (axis cs:90,-3.0) {top/bottom $C_p=-3$};
\filldraw[engred] (axis cs:90,-3) circle (1.6pt);
\node[engblack, font=\scriptsize, anchor=south] at (axis cs:135,0.05) {$C_p>0$ recovers at rear};
\end{axis}
\end{tikzpicture}
\caption[Inviscid pressure coefficient on a circular
cylinder]{Surface pressure coefficient on a circular cylinder in
steady inviscid uniform flow, $C_p = 1 - 4\sin^2\theta$, plotted from
the front stagnation point ($\theta=0$) to the rear stagnation point
($\theta=180^\circ$). The pressure peaks at both stagnation points
($C_p=1$) and reaches its minimum at the top and bottom
($\theta=90^\circ$, $C_p=-3$), where the surface speed is twice the
freestream. The distribution is symmetric front to back, so the
pressure pushing on the front is exactly recovered at the rear and
the net streamwise force integrates to zero: d'Alembert's paradox.
Real flow separates behind the cylinder, the rear pressure never
recovers, and the cylinder feels drag. Computed from the analytical
inviscid solution.}
\label{fig:vol03ch11:cylinder-cp}
\end{figure}


Figure~\ref{fig:vol03ch11:cylinder-cp} plots $C_p = 1 - 4\sin^2\theta$
around the surface. The front-to-back symmetry of the curve is the
whole of the paradox: whatever the front pressure pushes against, the
rear pressure pushes back equally, and the streamwise integral
cancels. The script \path{code/cylinder_potential.py} integrates the
surface pressure numerically and recovers a drag indistinguishable
from zero (about $10^{-15}\,\si{\newton\per\meter}$, machine round-off)
for the non-lifting case, and a lift equal to $\rho U_\infty\Gamma$ to
the same precision once a circulation is added; the surface
distribution it writes to \path{data/cylinder_cp.csv} reproduces the
analytical curve. The paradox is not a defect of the arithmetic; it
is a true property of the inviscid model, and its resolution waits for
the boundary layer of chapter 12.

\begin{mastery}
\textbf{The complex potential and conformal mapping.} For
two-dimensional irrotational incompressible flow the velocity
potential $\phi$ and stream function $\psi$ are harmonic conjugates,
so they combine into a single analytic function of the complex
coordinate $z = x + iy$, the \engterm{complex potential}
\[
F(z) = \phi(x,y) + i\,\psi(x,y).
\]
Because $F$ is analytic, its derivative is the \engterm{complex
velocity}
\[
\frac{dF}{dz} = u_x - i\,u_y,
\]
a single differentiation that returns both velocity components. The
elementary flows take compact closed forms: a uniform stream is
$F = U_\infty z$, a source of strength $m$ is $F = (m/2\pi)\ln z$, a
vortex of circulation $\Gamma$ is $F = -i(\Gamma/2\pi)\ln z$, and a
doublet is $F = -\kappa/(2\pi z)$. The lifting cylinder is their sum,
\[
F(z) = U_\infty\!\left(z + \frac{R^2}{z}\right)
- \frac{i\Gamma}{2\pi}\ln z .
\]
The power of the complex form is \engterm{conformal mapping}: an
analytic map $\zeta = g(z)$ deforms one flow domain into another while
preserving Laplace's equation, so a flow solved in a simple geometry
transfers to a complicated one for the price of an algebraic
substitution. The \engterm{Joukowski transform}
$\zeta = z + c^2/z$ carries the circle of the lifting cylinder into a
family of airfoil-shaped contours, turning the exactly solvable
cylinder-with-circulation into an exactly solvable cambered airfoil,
the first analytic airfoil theory (Joukowski, 1910). The same
machinery solves flow around corners, through slots, and over steps;
it is the analytic counterpart of the panel methods the simulation
exercises build numerically.
\end{mastery}

\subsection{The cylinder with circulation}

Adding a point vortex of circulation $\Gamma$ at the centre of the
cylinder produces an asymmetric flow with a net transverse force.
The stream function becomes
\[
\psi = U_\infty r \sin\theta\left(1 - \frac{R^2}{r^2}\right)
- \frac{\Gamma}{2\pi}\ln\frac{r}{R}.
\]
The surface pressure distribution is no longer symmetric about
$\theta = 0$ to $\pi$: integration gives a transverse force per
unit length
\[
L' = \rho U_\infty \Gamma,
\]
the \engterm{Kutta-Joukowski theorem} (Martin Kutta, 1902; Nikolai
Joukowski, 1906). The lift per unit length of any two-dimensional
body in an inviscid uniform flow equals the product of the fluid
density, the freestream speed, and the bound circulation around the
body. The theorem is the analytical foundation of airfoil theory
and the working basis of subsonic aerodynamic design.

\begin{mastery}
\textbf{The method of images: walls and free surfaces for free.} A
boundary that the flow cannot cross is a streamline, and the
linearity of Laplace's equation lets us build that streamline by
placing a mirror singularity outside the flow domain rather than
solving a boundary-value problem. The technique is borrowed intact
from electrostatics, where a grounded conducting plane is reproduced
by an image charge. Consider a point source of strength $m$ at height
$d$ above a flat solid wall along $y = 0$. The wall must carry zero
normal velocity. Placing an identical image source of strength $m$ at
$y = -d$, the mirror position, makes the two sources' vertical
velocities cancel exactly on $y = 0$ by symmetry, so the wall plane
is a streamline and the flow in the upper half-plane is the
superposition of the real source and its image. A vortex near a wall
takes an image vortex of opposite sign (so the tangential velocities
add and the normal velocities cancel), and the real vortex then
drifts parallel to the wall under the velocity its image induces, the
mechanism by which a wingtip vortex tracks along the ground during
takeoff and landing and by which a vortex ring approaching a wall
slows and spreads. A free surface, where the pressure rather than the
normal velocity is fixed, takes an image of the opposite kind: a
source images to a sink, so that the combined flow leaves the surface
undisturbed to first order. The same construction with a sequence of
images handles a source between two parallel walls (an infinite
image train) and a vortex inside a corner (a finite set of images
whose count depends on the corner angle). The method of images turns
a wall or a free surface into an algebra problem and is the first
tool an inviscid analyst reaches for whenever a boundary is flat.
\end{mastery}

\begin{archetype}[Balance, transport: Bernoulli plus Kutta-Joukowski]
The inviscid flow archetype puts two relations to work together.
Continuity sets the relationship between cross-sectional area and
flow speed; Bernoulli converts the speed variation into a pressure
variation; integrating the pressure over a body's surface gives the
net force. For a closed body in steady irrotational flow the
streamwise integration yields zero (d'Alembert's paradox); the
transverse integration yields the Kutta-Joukowski lift if the body
carries bound circulation. The archetype is the working basis of
every subsonic aerodynamics calculation, and it underwrites the
estimation habit of computing lift from a circulation estimate
rather than from an integration of surface pressures.
\end{archetype}

\begin{estimation}
\textbf{Estimate the pressure difference between the upper and
lower surfaces of an airliner wing at cruise, and compare it to
atmospheric pressure.}

\smallskip
\textit{Estimate, before reading on.}

\smallskip
A working airliner cruises at altitude near $10\,\si{\kilo\meter}$
where the ambient pressure is approximately $26\,\text{kPa}$ and
the air density is about $0.41\,\si{\kilogram\per\meter\cubed}$.
The cruise speed is about $250\,\si{\meter\per\second}$ (Mach
$0.85$). A medium twin-aisle airliner has a maximum takeoff weight
of about $230\,\text{t}$ ($2.3\times 10^6\,\si{\newton}$ of weight)
and a wing area of about $360\,\si{\meter}^2$.

The wing must supply lift equal to the aircraft's weight in cruise:
$L = 2.3\times 10^6\,\si{\newton}$. The average lift per unit area
is then
\[
\frac{L}{A_w} = \frac{2.3\times 10^6}{360} \approx 6.4\,\text{kPa}.
\]
The pressure differential between upper and lower surfaces,
averaged over the wing area, is about $6.4\,\text{kPa}$, roughly
a quarter of the ambient pressure at altitude.

The dynamic pressure at cruise is
\[
q = \tfrac{1}{2}\rho U_\infty^2 \approx \tfrac{1}{2}\times 0.41
\times 250^2 \approx 12.8\,\text{kPa}.
\]
The average pressure-coefficient differential between upper and
lower surfaces is then $\Delta C_p \approx 6.4/12.8 = 0.5$. The
working airfoil lift coefficient at cruise is typically
$C_L \approx 0.4$ to $0.6$, consistent with this estimate (the
chord-averaged $\Delta C_p$ over the wing area equals $C_L$ by
definition for a thin airfoil).

The postmortem: the rough match between $\Delta C_p \approx 0.5$
and the textbook $C_L$ at cruise shows that the simple force
balance ``lift equals weight'' plus the Bernoulli definition of
$C_p$ gives a quantitatively reasonable picture of cruise flight.
The pressure differential is not small relative to ambient: a
quarter of an atmosphere across a wing's chord is the working
mechanical reality. The next chapter sets a viscous correction
that costs a few percent of cruise thrust; this chapter's inviscid
analysis recovers the lift to leading order.
\end{estimation}

\section{Worked examples}\pathtag{standard}

We work three problems in detail. The reader who follows the
bookkeeping on these three should be able to apply Bernoulli plus
continuity to any standard inviscid configuration.

\subsection{Flow over an inviscid airfoil}

A symmetric thin airfoil of chord $c$ is placed at angle of attack
$\alpha$ in a uniform flow $U_\infty$. Thin-airfoil theory (Hermann
Glauert, 1926) approximates the airfoil as a line of bound vortices
distributed along the chord with circulation density $\gamma(x)$,
satisfying the boundary condition that the flow is tangent to the
camber line at every chord position. For a flat-plate symmetric
airfoil at small angle of attack, the result is
\[
C_L = 2\pi\alpha,
\]
the \engterm{thin-airfoil lift slope}. The lift coefficient is
linear in angle of attack with a slope of $2\pi$ per radian, or
about $0.11$ per degree.

The corresponding bound circulation per unit span is
\[
\Gamma = \pi c U_\infty \alpha,
\]
giving lift per unit span
\[
L' = \rho U_\infty \Gamma = \pi \rho U_\infty^2 c \alpha,
\]
in agreement with $L' = C_L \cdot \tfrac{1}{2}\rho U_\infty^2 c$
when $C_L = 2\pi\alpha$. The thin-airfoil result is the first
quantitative aerodynamic-design tool; real airfoils with finite
thickness and camber depart from it by perhaps $10$ to $20\,\%$ in
lift slope, more in the post-stall regime where the inviscid model
fails entirely (the topic of chapter 12 of this volume).

% Figure: airfoil with bound circulation; uniform stream plus circulation gives
% faster flow over the top, slower below, hence lift. Kutta condition at TE.
\begin{figure}[htbp]
\centering
\begin{tikzpicture}[
  foil/.style={draw=engblack, fill=enggray!25, very thick},
  upper/.style={draw=engred, -{Stealth}, very thick},
  lower/.style={draw=engblue, -{Stealth}, thick},
  circ/.style={draw=englime, -{Stealth}, very thick},
  label/.style={font=\scriptsize},
]
% Airfoil: a cambered shape via two bezier curves, chord ~5 units, slight AoA.
\draw[foil] (-2.5,-0.15)
  .. controls (-1.5,0.7) and (1.0,0.85) .. (2.5,0.05)   % upper surface
  .. controls (1.0,-0.25) and (-1.0,-0.45) .. (-2.5,-0.15); % lower surface
% Leading-edge stagnation marker
\filldraw[engblack] (-2.5,-0.15) circle (1.2pt);
% Trailing-edge: Kutta condition (smooth flow off the sharp TE)
\filldraw[engblack] (2.5,0.05) circle (1.2pt);
\node[label, anchor=west] at (2.6,0.05) {Kutta condition at TE};

% Upper-surface streamline (faster, longer arrow)
\draw[upper] (-3.6,1.0) .. controls (-1.5,1.4) and (1.0,1.4) .. (3.8,0.9);
\node[label, engred, anchor=south] at (0,1.45) {faster flow $\Rightarrow$ lower $p$};
% Lower-surface streamline (slower)
\draw[lower] (-3.6,-0.9) .. controls (-1.0,-1.0) and (1.0,-1.0) .. (3.8,-0.85);
\node[label, engblue, anchor=north] at (0,-1.0) {slower flow $\Rightarrow$ higher $p$};

% Bound circulation loop (green) around the foil
\draw[circ] (1.6,0.55) arc[start angle=0, end angle=300, x radius=1.6, y radius=0.95];
\node[label, englime, anchor=west] at (1.7,0.0) {$\Gamma$};

% Lift arrow
\draw[-{Stealth}, engamber, very thick] (0,0.3) -- (0,2.1);
\node[label, engamber, anchor=west] at (0.05,1.9) {$L' = \rho U_\infty \Gamma$};

% Freestream
\draw[-{Stealth}, engblack, thick] (-4.6,0) -- (-3.9,0);
\node[label, anchor=east] at (-4.6,0) {$U_\infty$};

\end{tikzpicture}
\caption[Bound circulation and lift on an airfoil]{The
Kutta-Joukowski picture of airfoil lift. The freestream $U_\infty$
plus a bound circulation $\Gamma$ (green loop) sum to a faster flow
over the upper surface (red) and a slower flow over the lower surface
(blue). By Bernoulli the faster upper flow carries a lower pressure
and the slower lower flow a higher pressure; the integrated pressure
difference is the lift $L' = \rho U_\infty \Gamma$ per unit span. The
circulation is not free to take any value: the Kutta condition fixes
it at exactly the value that lets the flow leave the sharp trailing
edge smoothly, which is how a real (slightly viscous) airfoil selects
its lift. The equal-transit-time story plays no part; particles over
the top do not meet particles from below at the trailing edge.}
\label{fig:vol03ch11:airfoil-circulation}
\end{figure}


Figure~\ref{fig:vol03ch11:airfoil-circulation} draws the
circulation picture the thin-airfoil result encodes: the freestream
and the bound circulation sum to a fast upper flow and a slow lower
flow, the pressure difference follows by Bernoulli, and the integrated
difference is the lift $\rho U_\infty\Gamma$. The Kutta condition is
the missing ingredient that fixes $\Gamma$; without it, the inviscid
equations admit any circulation and predict any lift.

\paragraph{Worked example: the lift slope as a sanity check.} A
flat-plate airfoil at $\alpha = 5^\circ = 0.0873\,\text{rad}$ in a
$60\,\si{\meter\per\second}$ stream of density
$1.2\,\si{\kilogram\per\meter\cubed}$ and chord $1\,\si{\meter}$ has
$C_L = 2\pi\alpha = 0.548$. The lift per unit span is $L' = C_L\cdot
\tfrac12\rho U_\infty^2 c = 0.548\times 0.5\times 1.2\times 60^2\times
1 = 1184\,\si{\newton\per\meter}$, and the bound circulation is
$\Gamma = L'/(\rho U_\infty) = 1184/(1.2\times 60) =
16.4\,\si{\meter\squared\per\second}$. The same $\Gamma$ follows
directly from $\Gamma = \pi c U_\infty\alpha = \pi\times 1\times
60\times 0.0873 = 16.4\,\si{\meter\squared\per\second}$, confirming
that the lift-slope and circulation routes agree. The two independent
routes meeting is the check; a calculation in which they disagree has
an arithmetic error or a hidden inconsistency in the assumptions.

\subsection{Nozzle flow}

A converging nozzle reduces the cross-sectional area of a steady
incompressible flow from $A_1$ to $A_2$, accelerating the flow from
speed $u_1$ to $u_2 = u_1(A_1/A_2)$ and decreasing the pressure
from $p_1$ to $p_2 = p_1 - \tfrac{1}{2}\rho(u_2^2 - u_1^2)$.

For a converging-diverging nozzle the flow accelerates in the
converging section, reaches its maximum speed at the throat (the
minimum area), and decelerates in the diverging section, recovering
some of the pressure. In an inviscid analysis the pressure recovery
is complete: the exit pressure equals the inlet pressure if the
exit area equals the inlet area. In real flow, viscous and
separation losses make the recovery incomplete; a well-designed
diffuser recovers about $70$ to $90\,\%$ of the dynamic pressure
expended in the contraction. The discipline of diffuser design is
to keep the divergence angle below the separation threshold (about
$7^{\circ}$ on each side for unaided flow, more with vortex
generators or trips).

For a compressible flow the analysis changes qualitatively at the
throat: above a critical pressure ratio, the throat \engterm{
chokes} and the flow speed at the throat equals the local speed of
sound. Beyond the throat, a converging-diverging nozzle can
accelerate the flow supersonically. Chapter 13 of this volume
develops the compressible analysis; the incompressible Bernoulli
result of this chapter is accurate at Mach numbers below about
$0.3$.

\subsection{Dam spillway}

Water flows over the crest of a dam spillway and accelerates down
the spillway face. Applying Bernoulli from the reservoir (where the
water is nearly at rest at elevation $z_1$) to a point on the
spillway face at elevation $z_2$, with both points at atmospheric
pressure,
\[
\tfrac{1}{2} u_2^2 = g(z_1 - z_2),
\]
giving $u_2 = \sqrt{2g(z_1 - z_2)}$. The spillway accelerates the
water through a head equal to the difference in elevation; this is
Torricelli's law applied to a free-surface flow.

The discharge per unit width of crest is set by the crest geometry
and the head above the crest. For a sharp-crested rectangular weir
of width $b$ and head $H$ above the crest,
\[
Q = C_d \tfrac{2}{3} b\sqrt{2g}\,H^{3/2},
\]
where $C_d \approx 0.62$ is a discharge coefficient that accounts
for contraction of the jet at the crest. The $H^{3/2}$ scaling
comes from Bernoulli (giving the velocity scaling $\sqrt{H}$)
combined with the cross-sectional area scaling proportional to
$H$.

Dam-engineering practice uses extensive empirical tables of $C_d$
for various spillway geometries (ogee crests, broad-crested weirs,
side channels, chute spillways); the underlying physics is
Bernoulli plus continuity, with empirical corrections for
viscosity, surface roughness, and separation effects. A spillway
designed for a $1000$-year flood will be sized using these
relations; the design uncertainty is mainly in the flood event,
not in the Bernoulli analysis.

\subsection{Draining-tank time as an unsteady Bernoulli problem}

Torricelli gives the instantaneous efflux speed; the time to drain a
tank turns that instantaneous law into an ordinary differential
equation, and the integration is worth carrying through because the
square-root-of-head law produces a drain time that surprises the
intuition built on constant-rate emptying.

% Figure: draining-tank geometry and the square-root-of-head decay of the
% surface as it empties through a bottom orifice.
\begin{figure}[htbp]
\centering
\begin{tikzpicture}
% ==== Left: schematic of the tank ====
\begin{scope}[xshift=-4.2cm]
\draw[draw=engblack, very thick] (-1.4,3) -- (-1.4,0) -- (1.4,0) -- (1.4,3);
\fill[engblue!15] (-1.4,0) rectangle (1.4,2.2);
\draw[engblue, thick] (-1.4,2.2) -- (1.4,2.2);
% surface arrow falling
\draw[engred, -{Stealth}, thick] (0,2.2) -- (0,1.4);
\node[font=\scriptsize, engblue, anchor=west] at (-1.35,2.4) {free surface, area $A_s$};
% head dimension
\draw[enggray, {Stealth}-{Stealth}, thin] (1.0,0) -- (1.0,2.2);
\node[font=\scriptsize, anchor=west] at (1.05,1.1) {$h$};
% orifice and jet
\draw[engblack, very thick] (-0.25,0) -- (-0.25,-0.25) -- (0.25,-0.25) -- (0.25,0);
\draw[engblue, -{Stealth}, thick] (0,-0.25) -- (0,-1.0);
\node[font=\scriptsize, engblue, anchor=west] at (0.3,-0.7) {$u=\sqrt{2gh}$};
\node[font=\scriptsize, anchor=north] at (0,-0.3) {area $A_o$};
\end{scope}
% ==== Right: depth history h(t) ====
\begin{scope}[xshift=2.5cm]
\begin{axis}[
  width=8.4cm, height=6cm,
  xlabel={time $t$ (h)},
  ylabel={depth $h$ (m)},
  xmin=0, xmax=9.5, ymin=0, ymax=2.2,
  axis lines=left,
  grid=major, grid style={enggray!20},
  tick label style={font=\scriptsize},
  label style={font=\small},
]
% h(t) = (sqrt(h0) - (C_d A_o / A_s) sqrt(g/2) t)^2, empties at ~9 h.
% Parametrise with the empty time T=9 h: h(t)=h0 (1 - t/T)^2 for t<=T.
\addplot[engblue, very thick, samples=120, domain=0:9]
  {2.0*(1 - x/9)^2};
\node[engblue, font=\scriptsize, anchor=south west] at (axis cs:0.2,1.7)
  {$h(t)=h_0\,(1-t/T)^2$};
\node[engred, font=\scriptsize, anchor=south] at (axis cs:8.3,0.1) {empty at $T\approx9$ h};
\filldraw[engred] (axis cs:9,0) circle (1.6pt);
\end{axis}
\end{scope}
\end{tikzpicture}
\caption[Draining-tank geometry and depth history]{Left: a tank of
free-surface area $A_s$ draining through a bottom orifice of area
$A_o$. Torricelli sets the instantaneous jet speed $u=\sqrt{2gh}$ from
the current head $h$; volume conservation, $A_s\,dh/dt=-C_d A_o
\sqrt{2gh}$, makes the problem unsteady. Right: integrating that
balance gives a depth that falls as $h(t)=h_0\,(1-t/T)^2$, a parabola
in time reaching zero at the empty time $T=(A_s/C_d A_o)\sqrt{2h_0/g}$.
The square-root-of-head efflux law makes the last shallow layer drain
slowly. For the $25\,\si{\meter}\times10\,\si{\meter}$ pool
$2\,\si{\meter}$ deep with a $100\,\si{\milli\meter}$ outlet, $T$ is
about nine hours. Computed from the Torricelli efflux law and volume
conservation.}
\label{fig:vol03ch11:draining-tank}
\end{figure}


\paragraph{Worked example: draining a swimming pool.} A rectangular
pool of plan area $A_s = 25\times 10 = 250\,\si{\meter}^2$, filled to
depth $h_0 = 2\,\si{\meter}$, drains through a bottom orifice of
diameter $100\,\si{\milli\meter}$ ($A_o = \tfrac{\pi}{4}(0.1)^2 =
7.85\times 10^{-3}\,\si{\meter}^2$) with discharge coefficient $C_d =
0.62$ (Figure~\ref{fig:vol03ch11:draining-tank}). Volume conservation
equates the rate the surface drops to the rate volume leaves the
orifice:
\[
A_s\,\frac{dh}{dt} = -C_d A_o \sqrt{2 g h}.
\]
This is separable. Writing $\beta = C_d A_o \sqrt{2g}/A_s$,
\[
\frac{dh}{\sqrt{h}} = -\beta\,dt
\quad\Longrightarrow\quad
2\sqrt{h} = 2\sqrt{h_0} - \beta t,
\]
so the depth history is the parabola $h(t) = \big(\sqrt{h_0} -
\tfrac{\beta}{2}t\big)^2$, and the tank is empty ($h = 0$) at
\[
T = \frac{2\sqrt{h_0}}{\beta}
= \frac{A_s}{C_d A_o}\sqrt{\frac{2 h_0}{g}}.
\]
Putting in numbers, $\sqrt{2 h_0/g} = \sqrt{4/9.81} =
0.639\,\si{\second}$ (carrying $\si{\meter}^{1/2}$ implicitly), and
$A_s/(C_d A_o) = 250/(0.62\times 7.85\times 10^{-3}) = 5.14\times
10^4$, so
\[
T = 5.14\times 10^4 \times 0.639 = 3.28\times 10^4\,\si{\second}
\approx 9.1\,\text{h}.
\]
A pool that would empty in about three hours at the initial efflux
rate of $C_d A_o\sqrt{2 g h_0} = 0.62\times 7.85\times 10^{-3}\times
6.26 = 0.030\,\si{\meter\cubed\per\second}$ (which clears the
$500\,\si{\meter}^3$ in $1.6\times 10^4\,\si{\second}\approx
4.6\,\text{h}$) in fact takes twice as long, because the head, and
with it the efflux speed, falls as the tank empties. The script
\path{code/tank_drain.py} integrates the same ODE numerically and
recovers the nine-hour figure, a check that the closed form and the
numerics agree.

\subsection{Pitot-static airspeed and the compressibility correction}

The incompressible Pitot relation $u = \sqrt{2(p_0 - p)/\rho}$ is the
working tool of low-speed anemometry; at flight Mach numbers the same
measurement needs a correction, and seeing where the correction
enters is the point of the example.

\paragraph{Worked example: airspeed at low and at high subsonic
Mach.} A Pitot-static probe reads a stagnation-static differential of
$\Delta p = p_0 - p = 4.0\,\text{kPa}$ in air of static pressure $p =
70\,\text{kPa}$ and density $\rho = 0.91\,\si{\kilogram\per\meter
\cubed}$ (about $3\,\si{\kilo\meter}$ altitude). The incompressible
reading is
\[
u = \sqrt{\frac{2\Delta p}{\rho}} = \sqrt{\frac{2\times 4000}{0.91}}
= \sqrt{8791} = 93.8\,\si{\meter\per\second}.
\]
The local speed of sound is $c = \sqrt{\gamma p/\rho} = \sqrt{1.4
\times 70\,000/0.91} = \sqrt{1.077\times 10^5} =
328\,\si{\meter\per\second}$, so the Mach number is $M = u/c = 93.8/
328 = 0.29$, just at the edge of the incompressible regime. The
compressible stagnation relation for isentropic flow of a perfect gas
is
\[
\frac{p_0}{p} = \left(1 + \frac{\gamma - 1}{2}M^2\right)
^{\gamma/(\gamma - 1)},
\]
which at $M = 0.29$ gives $p_0/p = (1 + 0.2\times 0.0841)^{3.5} =
(1.0168)^{3.5} = 1.0601$, against the incompressible prediction
$1 + \tfrac12\gamma M^2 = 1 + 0.7\times 0.0841 = 1.0589$. The two
agree to about $0.1\,\%$ in pressure ratio at this Mach number, so
the incompressible airspeed is good to a fraction of a percent here.
Repeating at $M = 0.7$ (a transport cruise condition), the
compressible ratio is $(1 + 0.2\times 0.49)^{3.5} = (1.098)^{3.5} =
1.387$ while the incompressible form gives $1 + 0.7\times 0.49 =
1.343$, a $3\,\%$ discrepancy in pressure ratio that translates to a
few percent in computed speed. The lesson sets the boundary: below
$M \approx 0.3$ the incompressible Pitot relation is accurate to under
a percent; above it the compressible stagnation relation of
chapter 13 of this volume must replace it, and an airspeed indicator
that ignores the correction reads progressively low as Mach rises.

\subsection{Orifice metering and the cavitation margin}

A metering orifice sized for a clean differential-pressure signal can
inadvertently drive its own throat into cavitation; the example
closes the loop between the Venturi/orifice metering physics and the
cavitation failure mode by carrying both calculations on one device.

% Figure: orifice plate with the vena contracta downstream and the static
% pressure profile dipping at the contraction, marking the cavitation margin.
\begin{figure}[htbp]
\centering
\begin{tikzpicture}[
  wall/.style={draw=engblack, very thick},
  jet/.style={draw=engblue, thick},
  flow/.style={draw=engblue, -{Stealth}, thick},
  label/.style={font=\scriptsize},
]
% ==== pipe walls ====
\draw[wall] (-4,1.1) -- (4,1.1);
\draw[wall] (-4,-1.1) -- (4,-1.1);
% ==== orifice plate at x=0 ====
\draw[wall, fill=enggray!25] (-0.12,1.1) -- (0.12,1.1) -- (0.12,0.45) -- (-0.12,0.45) -- cycle;
\draw[wall, fill=enggray!25] (-0.12,-1.1) -- (0.12,-1.1) -- (0.12,-0.45) -- (-0.12,-0.45) -- cycle;
\node[label, anchor=south] at (0,1.15) {orifice plate};
% ==== streamlines through the hole, contracting to a vena contracta ====
\draw[jet] (-4,0.9) .. controls (-1.2,0.85) and (-0.3,0.45) .. (0.12,0.42)
  .. controls (0.7,0.36) and (1.0,0.30) .. (1.1,0.30)
  .. controls (1.6,0.32) and (2.6,0.55) .. (4,0.7);
\draw[jet] (-4,-0.9) .. controls (-1.2,-0.85) and (-0.3,-0.45) .. (0.12,-0.42)
  .. controls (0.7,-0.36) and (1.0,-0.30) .. (1.1,-0.30)
  .. controls (1.6,-0.32) and (2.6,-0.55) .. (4,-0.7);
\draw[flow] (-3.2,0) -- (-2.0,0);
% vena contracta marker
\draw[engred, {Stealth}-{Stealth}, thin] (1.1,0.30) -- (1.1,-0.30);
\node[label, engred, anchor=west] at (1.2,0) {vena contracta $A_c$};
\node[label, anchor=north] at (0,-1.15) {area $A_2$};
\node[label, anchor=north] at (-2.6,-1.15) {area $A_1$, $p_1$};
% ==== static-pressure mini-plot below ====
\begin{scope}[yshift=-3.6cm]
\draw[enggray, -{Stealth}] (-4,0) -- (4.2,0) node[label, anchor=west] {$x$};
\draw[enggray, -{Stealth}] (-4,-0.2) -- (-4,1.8) node[label, anchor=south] {$p$};
% pressure: high upstream, deep dip at vena contracta, partial recovery
\draw[engblue, very thick]
  (-4,1.5) .. controls (-1.5,1.45) and (-0.5,1.0) .. (1.1,0.15)
  .. controls (2.0,0.6) and (3.0,0.95) .. (4,1.0);
\node[label, engblue, anchor=south west] at (-3.9,1.5) {$p_1$};
\node[label, engblue, anchor=north] at (1.1,0.1) {$p_{\min}$};
% vapour pressure margin line
\draw[engred, dashed, very thick] (-4,0.32) -- (4,0.32);
\node[label, engred, anchor=west] at (3.0,0.32) {$p_v$};
\draw[engamber, {Stealth}-{Stealth}, thin] (1.1,0.15) -- (1.1,0.32);
\node[label, engamber, anchor=west] at (1.2,0.24) {margin};
\end{scope}
\end{tikzpicture}
\caption[Orifice-plate flow and the cavitation margin]{Top: flow
through a thin orifice plate. The streamlines continue to contract
past the plate to a minimum-area \emph{vena contracta} of area
$A_c<A_2$, where the speed is highest and the static pressure lowest;
the discharge coefficient $C_d\approx0.6$ folds the contraction and a
small viscous loss into one factor. Bottom: the static pressure along
the centreline, high upstream, dipping to $p_{\min}$ at the vena
contracta, and recovering only partially downstream. The cavitation
margin is the gap between $p_{\min}$ and the vapour pressure $p_v$; a
metering orifice sized for a large pressure drop can push $p_{\min}$
below $p_v$ and cavitate, so the design discipline is to check the
margin at the worked flow. Computed from continuity, Bernoulli, and
the vena-contracta contraction.}
\label{fig:vol03ch11:orifice-vena-contracta}
\end{figure}


\paragraph{Worked example: an orifice plate near its cavitation
limit.} An orifice plate of bore diameter $40\,\si{\milli\meter}$
meters water in a $100\,\si{\milli\meter}$ line at $Q =
12\,\text{L/s}$, with the line at a static pressure of $p_1 =
180\,\text{kPa}$ absolute (Figure~\ref{fig:vol03ch11:orifice-vena-contracta}).
The bore area is $A_2 = \tfrac{\pi}{4}(0.04)^2 = 1.257\times
10^{-3}\,\si{\meter}^2$, and the flow contracts further to a vena
contracta of area $A_c \approx C_c A_2$ with contraction coefficient
$C_c \approx 0.62$, so $A_c = 7.79\times 10^{-4}\,\si{\meter}^2$. The
speed at the vena contracta is
\[
u_c = \frac{Q}{A_c} = \frac{0.012}{7.79\times 10^{-4}}
= 15.4\,\si{\meter\per\second}.
\]
The approach speed in the line is $u_1 = Q/A_1 = 0.012/(7.85\times
10^{-3}) = 1.53\,\si{\meter\per\second}$, negligible by comparison.
Bernoulli from the line to the vena contracta gives the minimum
pressure
\[
p_\text{min} = p_1 - \tfrac12\rho(u_c^2 - u_1^2)
= 180\,000 - 0.5\times 1000\times(237 - 2.3)
= 180\,000 - 117\,400 = 62.6\,\text{kPa}.
\]
The cavitation margin is the gap to the vapour pressure $p_v \approx
2.3\,\text{kPa}$ absolute, expressed as the cavitation number
\[
\sigma = \frac{p_\text{min} - p_v}{\tfrac12\rho u_c^2}
= \frac{62\,600 - 2300}{117\,400} = 0.51.
\]
For sharp-edged orifices the inception value of $\sigma$ is typically
near $0.6$ to $1.5$ depending on the recovery geometry, so $\sigma =
0.51$ sits inside the cavitating range: this orifice cavitates at this
flow, and the downstream face and pipe wall will pit over time. The
fix is to lower the velocity at the contraction by opening the bore
(a $50\,\si{\milli\meter}$ bore drops $u_c$ to about
$9.9\,\si{\meter\per\second}$ and $\tfrac12\rho u_c^2$ to
$49\,\text{kPa}$, lifting $\sigma$ above $2.5$), at the cost of a
smaller and noisier differential signal, or to raise the line
pressure with downstream backpressure. The metering accuracy and the
cavitation margin pull against each other, and the design must satisfy
both.

\begin{estimation}
\textbf{Estimate the velocity and the cavitation risk at the foot of
a $30\,\si{\meter}$ dam spillway.}

\smallskip
\textit{Estimate, before reading on.}

\smallskip
Water released at the crest is nearly at rest; at the foot it has
fallen a head of about $30\,\si{\meter}$. Torricelli on the
free-surface streamline gives $u = \sqrt{2gh} = \sqrt{2\times 9.81
\times 30} \approx 24\,\si{\meter\per\second}$. At that speed the
dynamic pressure is $q = \tfrac12\rho u^2 = 0.5\times 1000\times 24^2
\approx 290\,\text{kPa}$.

Where the high-speed sheet passes over a surface irregularity (a
construction joint, an offset, a slot for a gate), the local static
pressure can drop by a sizeable fraction of $q$ below the ambient
atmospheric value of about $100\,\text{kPa}$ absolute. A local
$C_p$ of only $-0.4$ at that speed gives a pressure drop of
$0.4\times 290 = 116\,\text{kPa}$, pulling the local absolute pressure
to roughly $100 - 116 < 0$, below the vapour pressure of about
$2.3\,\text{kPa}$: cavitation. The cavitation number there is
$\sigma = (100 - 2.3)/290 \approx 0.34$, near the inception range for
surface irregularities.

\smallskip
The postmortem: the estimate explains why high-head spillways are
aerated. Introducing air at the surface raises the local pressure and
cushions bubble collapse; the design discipline is to keep $\sigma$
above the inception value at every irregularity, or to aerate where
that cannot be guaranteed. The Bernoulli pressure drop is the
quantity that flags the hazard, and a few tens of metres of head is
enough to reach it.
\end{estimation}

\section{Failure: where Bernoulli does not apply, and why students misuse it}\pathtag{core}

Every chapter in this work has a failure section. The failure mode
for inviscid analysis is qualitatively different from the failure
modes of structural mechanics: an inviscid Bernoulli calculation
applied to a configuration that violates its assumptions gives a
plausible-looking number that is wrong, sometimes by an order of
magnitude.

\subsection{Bernoulli across a hydraulic jump}

Water flowing rapidly down a spillway and entering a slower
tailrace transitions through a \engterm{hydraulic jump}: an abrupt
rise in water depth from a supercritical fast-shallow flow to a
subcritical slow-deep flow. Bernoulli applied across the jump
predicts a pressure rise consistent with the velocity decrease,
but the actual energy loss is substantial: the jump dissipates a
fraction of the upstream kinetic energy as turbulence, heat, and
acoustic radiation. The fractional energy loss can exceed $50\,\%$
for strong jumps (Froude numbers above $5$).

The reason Bernoulli fails across the jump is that the flow is not
inviscid in the jump region: intense turbulent dissipation
converts mechanical energy to heat. The control-volume momentum
balance still holds (it makes no inviscid assumption), and the
correct analytical tool is the momentum equation across the jump.
The momentum analysis gives the downstream depth and the energy
loss; Bernoulli applied through the jump gives the wrong depth
and conserves an energy that is in fact being dissipated.

\subsection{Bernoulli inside a boundary layer}

A real flow past a solid surface develops a thin boundary layer
(chapter 12 of this volume) in which viscous stresses dominate.
Inside the boundary layer, the streamwise velocity drops from the
freestream value to zero at the wall. Bernoulli predicts that the
pressure must rise correspondingly along the streamline that
approaches the wall, by an amount $\tfrac{1}{2}\rho U_\infty^2$.

The actual pressure distribution at the wall is set by the
inviscid flow just outside the boundary layer, not by the
within-boundary-layer streamline. The boundary-layer approximation
(Prandtl, 1904; chapter 12 of this volume) states that the
pressure is constant across the boundary layer (in the wall-normal
direction): the boundary layer is so thin that the pressure
imposed by the outer inviscid flow penetrates to the wall
unchanged. A Bernoulli calculation along a streamline that crosses
into the boundary layer is therefore qualitatively wrong: the
pressure does not rise to the stagnation value, and the velocity
loss is due to viscous dissipation, not pressure recovery.

\subsection{The lift formula misused}

A common misconception holds that an airfoil generates lift
because air takes longer to travel over the longer-curved upper
surface than over the shorter lower surface, and ``must''
therefore travel faster on the upper surface to ``meet up'' with
the lower-surface flow at the trailing edge. The argument is the
\engterm{equal-transit-time fallacy}, and it is wrong on every
point: there is no physical principle that forces fluid particles
to meet up at the trailing edge (they do not), the upper-surface
flow on a real airfoil is much faster than the equal-transit
calculation predicts (often by a factor of two or more), and flat
plates and inverted cambered airfoils generate lift in
configurations where the equal-transit argument predicts no lift
or negative lift.

The correct explanation rests on the Kutta-Joukowski theorem and
the establishment of bound circulation through the Kutta
condition: a real airfoil at angle of attack develops a sharp
trailing-edge flow pattern that fixes the circulation at the value
required to keep the rear stagnation point at the trailing edge
itself (the \engterm{Kutta condition}, Martin Kutta, 1902). The
bound circulation, combined with the freestream, produces a
faster flow on the upper surface and a slower flow on the lower
surface; the pressure differential follows by Bernoulli; the
integrated pressure differential is the lift, equal to
$\rho U_\infty \Gamma$. The argument is closed once the
circulation is fixed by the Kutta condition.

\subsection{Bernoulli in unsteady or rotational flow}

A pulsating flow, a starting transient, or a flow with significant
vorticity along the streamline of interest all violate Bernoulli's
assumptions. The correct generalisation for unsteady irrotational
flow is the \engterm{unsteady Bernoulli equation}
\[
\frac{\partial \phi}{\partial t} + \tfrac{1}{2}u^2 + \frac{p}{\rho}
+ gz = f(t),
\]
which carries an additional time-derivative term. For a rotational
flow, the Bernoulli constant differs between streamlines, and the
relation cannot be applied across streamlines without an analysis
of the vorticity distribution. Students who memorise the
incompressible steady form and apply it to oscillating pipe flows,
starting jets, or vortex-shedding wakes get numerically incorrect
answers; the discipline is to check the five assumptions before
writing the equation.

\subsection{Cavitation: when Bernoulli predicts a pressure the
liquid cannot sustain}

A liquid cannot carry an absolute pressure below its vapour pressure
$p_v$; at that pressure it boils. Bernoulli is the calculation that
warns when a design approaches this limit. Where a flow accelerates
through a constriction, a pump inlet, a valve, a propeller-blade
suction surface, or the toe of a spillway, the static pressure drops
by $\tfrac12\rho(u_2^2 - u_1^2)$, and a high enough local speed drives
the absolute pressure below $p_v$. Vapour bubbles form, are swept
downstream into a higher-pressure region, and collapse. The collapse
is violent: the bubble walls accelerate inward and the implosion sends
a pressure spike of hundreds of megapascals into any nearby surface,
pitting bronze impellers, steel valve seats, and concrete spillway
faces over time. The phenomenon is \engterm{cavitation}, and its onset
is governed not by viscosity but by the inviscid Bernoulli pressure
drop.

% Figure: pressure along a streamline through a constriction dropping below
% vapour pressure, marking the cavitation zone.
\begin{figure}[htbp]
\centering
\begin{tikzpicture}
\begin{axis}[
  width=12cm, height=7cm,
  xlabel={position along the flow $s$ (arb.\ units)},
  ylabel={static pressure $p$ (kPa, absolute)},
  xmin=0, xmax=10, ymin=-5, ymax=120,
  axis lines=left,
  grid=major, grid style={enggray!20},
  tick label style={font=\scriptsize},
  label style={font=\small},
  legend style={font=\scriptsize},
  legend pos=north east,
]
% Static pressure: high upstream, dips deeply at throat (s=5), partial recovery.
\addplot[engblue, very thick, samples=200, domain=0:10]
  {95 - 110*exp(-((x-5)^2)/1.3)};
\addlegendentry{static pressure $p(s)$}
% Vapour pressure of water at ~20 C (about 2.3 kPa absolute)
\addplot[engred, dashed, very thick, samples=2, domain=0:10] {2.3};
\addlegendentry{vapour pressure $p_v\approx2.3$ kPa}
% Cavitation region shading where p < p_v: roughly s in [4.4, 5.6]
\fill[engred!18] (axis cs:4.35,-5) rectangle (axis cs:5.65,2.3);
\node[engred, font=\scriptsize, anchor=north] at (axis cs:5,18) {cavitation zone};
\node[engred, font=\scriptsize, anchor=south] at (axis cs:5,2.3) {$p<p_v$: vapour bubbles form};
\end{axis}
\end{tikzpicture}
\caption[Cavitation onset along a streamline]{Static pressure along a
streamline through a constriction (a valve, a pump inlet, or a
spillway aerator step). Bernoulli converts the high throat velocity
into a deep pressure drop; where the local static pressure falls below
the liquid's vapour pressure $p_v$ (red dashed line, about
$2.3\,\text{kPa}$ absolute for water near $20\,^{\circ}\mathrm{C}$),
the water boils locally and vapour bubbles form. Downstream, where the
pressure recovers, the bubbles collapse violently and the implosion
pressure pits metal and concrete: cavitation damage. The inviscid
Bernoulli pressure prediction is exactly the calculation that flags
cavitation risk before it happens. Computed from a Gaussian throat
pressure model.}
\label{fig:vol03ch11:cavitation}
\end{figure}


Figure~\ref{fig:vol03ch11:cavitation} traces the static pressure
along a streamline through a constriction down past the vapour
pressure of water (about $2.3\,\text{kPa}$ absolute near
$20\,^{\circ}\mathrm{C}$). The dimensionless group that ranks
cavitation risk is the \engterm{cavitation number}
\[
\sigma = \frac{p - p_v}{\tfrac12\rho U^2},
\]
the margin between local static pressure and vapour pressure measured
in dynamic pressures; cavitation begins when $\sigma$ falls to a
geometry-dependent inception value. The lesson for the failure
section is that Bernoulli here gives the right answer for the wrong
comfort: the inviscid pressure prediction is accurate up to the point
of inception, and it is exactly that accurate prediction of a very low
pressure that flags the hazard. A designer who computes the throat
pressure and finds it below $p_v$ has been warned by Bernoulli; one
who never computes it learns about cavitation from the pitted hardware.

\subsection{Water hammer: when the steady assumption breaks
catastrophically}

Closing a valve at the end of a pipe carrying flowing liquid stops the
flow, and the moving column of liquid cannot stop instantly. The
kinetic energy of the column converts into a pressure surge that
travels back up the pipe as an acoustic wave at the speed of sound in
the confined liquid, reflects, and returns, ringing the pipe like a
struck bell. The phenomenon is \engterm{water hammer}, the audible
bang in household plumbing when a tap or appliance valve shuts
quickly. Steady Bernoulli has nothing to say about it: the flow is
violently unsteady, and the relevant term is the $\partial\phi/
\partial t$ contribution that steady Bernoulli discards.

The peak pressure rise for an instantaneous stop is the
\engterm{Joukowsky relation}
\[
\Delta p = \rho a\,\Delta u,
\]
where $a$ is the pressure-wave speed (close to the liquid's sonic
speed, reduced by pipe-wall elasticity) and $\Delta u$ is the change
in flow speed. For water with $a \approx 1200\,\si{\meter\per\second}$
and a flow brought to rest from only $2\,\si{\meter\per\second}$, the
surge is $\Delta p = 1000\times 1200\times 2 = 2.4\,\text{MPa}$, some
twenty-four atmospheres, far above the steady operating pressure and
enough to rupture fittings. The surge is bounded only if the valve
closes slowly relative to the wave's round-trip time $2L/a$ (with $L$
the pipe length); a slow closure lets successive reflections partly
cancel the surge. Surge tanks, air chambers, and pressure-relief
valves are the standard mitigations. Water hammer is the
unsteady-flow failure mode of this chapter made concrete: the steady
assumption is not approximately wrong, it discards the entire
physics.

\subsection{Case study: water-hammer surge in a hydroelectric
penstock}

The water-hammer failure mode deserves a full end-to-end treatment,
because the engineering decision it forces, how fast a valve may be
allowed to close, follows from the inviscid relations of this chapter
strung together with the elasticity of the pipe and the geometry of
the installation. We carry one penstock through from its steady
operating state to a verdict on its pressure rating.

% Figure: hydroelectric penstock layout. Reservoir, surge tank, sloping
% penstock, turbine valve at the foot. The geometry the water-hammer case
% study carries numbers through.
\begin{figure}[htbp]
\centering
\begin{tikzpicture}[
  wall/.style={draw=engblack, very thick},
  pipe/.style={draw=engblue, very thick},
  water/.style={fill=engblue!15},
  flow/.style={draw=engblue, -{Stealth}, thick},
  dim/.style={draw=enggray, {Stealth}-{Stealth}, thin},
  label/.style={font=\scriptsize},
]
% Reservoir block (top left)
\fill[water] (-5,3.2) rectangle (-2.2,4.4);
\draw[wall] (-5,4.4) -- (-5,3.2) -- (-2.2,3.2);
\draw[engblue, thick] (-5,4.4) -- (-2.2,4.4);
\node[label, engblue] at (-3.6,4.0) {reservoir};
\node[label] at (-3.6,4.62) {free surface};
% Headrace from reservoir to surge tank
\draw[pipe] (-2.2,3.5) -- (-0.6,3.5);
% Surge tank (vertical riser)
\draw[wall] (-0.7,3.2) -- (-0.7,5.2);
\draw[wall] (-0.5,3.2) -- (-0.5,5.2);
\fill[water] (-0.7,3.2) rectangle (-0.5,4.7);
\node[label, anchor=west] at (-0.45,5.0) {surge tank};
% Penstock: sloping pipe down to the valve
\draw[pipe] (-0.6,3.35) -- (4.0,-1.2);
\draw[pipe] (-0.6,3.05) -- (4.2,-1.45);
% Flow arrow along penstock
\draw[flow] (0.6,2.4) -- (2.2,0.85);
\node[label, engblue] at (1.8,1.9) {$u_0$};
% Valve at the foot
\draw[wall, fill=engamber!30] (3.95,-1.25) rectangle (4.35,-1.0);
\node[label, anchor=west] at (4.4,-1.1) {valve};
% Turbine downstream of valve
\draw[wall] (4.6,-1.6) -- (5.3,-1.6) -- (5.3,-0.7) -- (4.6,-0.7) -- cycle;
\node[label] at (4.95,-1.15) {turbine};
% Datum line at the foot
\draw[enggray, dashed] (-5.2,-1.5) -- (5.6,-1.5);
\node[label, enggray, anchor=west] at (5.0,-1.7) {datum};
% Gross-head dimension
\draw[dim] (-5.6,-1.5) -- (-5.6,4.4);
\node[label, anchor=east] at (-5.65,1.5) {gross head $H$};
% Penstock length label
\node[label, engblue, rotate=-44] at (2.0,1.2) {length $L$};
\end{tikzpicture}
\caption[Hydroelectric penstock layout for the water-hammer case
study]{Layout of a hydroelectric scheme: a reservoir feeds a sloping
penstock of length $L$ that drops a gross head $H$ to a valve and
turbine at the foot. A surge tank near the top of the penstock gives
the water column a free surface to vent into when the valve closes,
limiting the pressure rise. The water-hammer case study carries
numbers through this geometry: steady flow at speed $u_0$, sudden
valve closure, the Joukowsky pressure rise, and the surge-tank or
slow-closure mitigation that keeps the peak pressure inside the pipe's
rating.}
\label{fig:vol03ch11:penstock-layout}
\end{figure}


\paragraph{The installation.} A small hydroelectric scheme
(Figure~\ref{fig:vol03ch11:penstock-layout}) draws from a reservoir
through a steel penstock of length $L = 600\,\si{\meter}$, inside
diameter $D = 1.2\,\si{\meter}$, and wall thickness $e =
12\,\si{\milli\meter}$, dropping a gross head $H = 90\,\si{\meter}$ to
a turbine. The penstock steel has Young's modulus $E =
200\,\text{GPa}$; the water has density $\rho =
1000\,\si{\kilogram\per\meter\cubed}$ and bulk modulus $K =
2.2\,\text{GPa}$. At full load the turbine passes $Q =
3.4\,\si{\meter\cubed\per\second}$. The pipe is rated to a maximum
working pressure of $1.6\,\text{MPa}$. The question is whether a fast
emergency closure of the turbine inlet valve can be survived, and if
not, what closure schedule or mitigation keeps the surge inside the
rating.

\paragraph{Step 1: the steady state.} The penstock area is $A =
\tfrac{\pi}{4}D^2 = \tfrac{\pi}{4}(1.2)^2 = 1.131\,\si{\meter}^2$, so
the steady flow speed is
\[
u_0 = \frac{Q}{A} = \frac{3.4}{1.131} = 3.0\,\si{\meter\per\second}.
\]
The static pressure at the foot, from Bernoulli on the free-surface
streamline (reservoir at rest, the foot moving at $u_0$, both at the
same kinetic correction), is dominated by the gross head:
$p_\text{foot} \approx \rho g H = 1000\times 9.81\times 90 =
0.88\,\text{MPa}$, comfortably inside the $1.6\,\text{MPa}$ rating.
Steady operation is not the concern; the transient is.

\paragraph{Step 2: the wave speed with pipe elasticity.} The pressure
wave does not travel at the sonic speed of the water alone, because
the pipe wall stretches as the pressure rises and stores some of the
fluid that the compression would otherwise have to absorb. The
Korteweg formula corrects the water sonic speed
$a_w = \sqrt{K/\rho}$ for the wall compliance of a thin-walled pipe:
\[
a = \sqrt{\frac{K/\rho}{1 + (K D)/(E e)}}.
\]
The water-alone speed is $a_w = \sqrt{2.2\times 10^9/1000} =
1483\,\si{\meter\per\second}$. The compliance term is
\[
\frac{K D}{E e} = \frac{(2.2\times 10^9)(1.2)}
{(200\times 10^9)(0.012)} = \frac{2.64\times 10^9}{2.4\times 10^9}
= 1.10,
\]
so
\[
a = \frac{1483}{\sqrt{1 + 1.10}} = \frac{1483}{1.449}
= 1024\,\si{\meter\per\second}.
\]
The elastic pipe slows the wave by nearly a third relative to rigid
confinement. We round to $a = 1000\,\si{\meter\per\second}$ for the
estimates that follow.

\paragraph{Step 3: the Joukowsky surge for instantaneous closure.}
Stopping the column instantaneously converts its momentum into a
pressure spike given by the Joukowsky relation $\Delta p = \rho a
\Delta u$, with $\Delta u = u_0 = 3.0\,\si{\meter\per\second}$:
\[
\Delta p = 1000\times 1000\times 3.0 = 3.0\,\text{MPa}.
\]
Added to the steady $0.88\,\text{MPa}$, the peak pressure at the
valve reaches $3.9\,\text{MPa}$, well above the $1.6\,\text{MPa}$
rating. An instantaneous (or sufficiently fast) closure ruptures the
penstock. The verdict at this step is a fail, and the rest of the
case study is the mitigation.

\paragraph{Step 4: the closure-time threshold.} The full Joukowsky
surge develops only when the valve closes before the relief wave can
return from the reservoir. The wave travels to the upper end and back
in the round-trip time
\[
t_r = \frac{2L}{a} = \frac{2\times 600}{1000} = 1.2\,\si{\second}.
\]
A closure faster than $1.2\,\si{\second}$ is hydraulically
``instantaneous'' and produces the full $3.0\,\text{MPa}$ surge. A
closure slower than $t_r$ lets successive reflections cancel part of
the rise; the arithmetic-surge approximation estimates the reduced
peak as
\[
\Delta p_\text{slow} \approx \rho a u_0 \,\frac{2L/a}{t_c}
= \frac{\rho L u_0}{t_c}\cdot 2,
\]
inversely proportional to the closure time
(Figure~\ref{fig:vol03ch11:waterhammer-surge}). To hold the surge to
a margin below the rating, say to $\Delta p_\text{allow} =
0.7\,\text{MPa}$ on top of the steady $0.88\,\text{MPa}$ (keeping the
peak at $1.58\,\text{MPa}$, just inside the rating), the closure must
satisfy
\[
t_c \ge \frac{2L u_0}{\Delta p_\text{allow}/\rho}
= \frac{2\times 600\times 3.0}{0.7\times 10^6/1000}
= \frac{3600}{700} = 5.1\,\si{\second}.
\]
A closure spread over at least about five seconds keeps the surge
inside the rating; the turbine governor and the valve actuator are
specified to that minimum.

% Figure: pressure rise at the valve versus closure time, showing the flat
% Joukowsky plateau for rapid closure and the falling branch once the closure
% time exceeds the wave round-trip time 2L/a.
\begin{figure}[htbp]
\centering
\begin{tikzpicture}
\begin{axis}[
  width=12cm, height=7cm,
  xlabel={valve closure time $t_c$ (s)},
  ylabel={peak pressure rise $\Delta p$ (MPa)},
  xmin=0, xmax=12, ymin=0, ymax=2.8,
  axis lines=left,
  grid=major, grid style={enggray!20},
  tick label style={font=\scriptsize},
  label style={font=\small},
  legend style={font=\scriptsize},
  legend pos=north east,
]
% Round-trip time 2L/a marked; for L=600 m, a=1000 m/s, 2L/a = 1.2 s.
% Rapid closure (t_c < 2L/a): full Joukowsky surge, flat plateau at 2.4 MPa.
\addplot[engblue, very thick, samples=2, domain=0:1.2] {2.4};
% Slow closure (t_c > 2L/a): approximate arithmetic surge dp ~ Joukowsky*(2L/a)/t_c.
\addplot[engblue, very thick, samples=150, domain=1.2:12] {2.4*1.2/x};
\addlegendentry{peak surge $\Delta p(t_c)$}
% Round-trip time marker
\addplot[engred, dashed, very thick, samples=2] coordinates {(1.2,0) (1.2,2.4)};
\addlegendentry{round-trip $2L/a=1.2$ s}
% Pipe rating line
\addplot[engamber, dotted, very thick, samples=2, domain=0:12] {1.0};
\addlegendentry{pipe rating $1.0$ MPa}
\node[engblue, font=\scriptsize, anchor=south] at (axis cs:0.6,2.42) {Joukowsky plateau};
\node[engred, font=\scriptsize, anchor=west] at (axis cs:1.3,1.9) {rapid $\mid$ slow};
\node[engamber, font=\scriptsize, anchor=south] at (axis cs:8,1.02) {safe closure: $t_c\gtrsim2.9$ s};
\end{axis}
\end{tikzpicture}
\caption[Water-hammer surge versus valve closure time]{Peak pressure
rise at the valve as a function of closure time for a penstock of
length $L=600\,\si{\meter}$ carrying water at $u_0=3\,\si{\meter\per
\second}$ with a wave speed $a=1000\,\si{\meter\per\second}$. For any
closure faster than the wave round-trip time $2L/a=1.2\,\si{\second}$,
the column never feels the reflected relief wave and the full
Joukowsky surge $\Delta p=\rho a u_0=3.0\,\text{MPa}$ develops
(the flat plateau, here drawn at the worked value $2.4\,\text{MPa}$
for $u_0=2.4\,\si{\meter\per\second}$ effective). For slower closures
the surge falls roughly as $\Delta p\approx\rho a u_0\,(2L/a)/t_c$,
the arithmetic-surge estimate, because successive reflections cancel
part of the rise. Holding the surge below the pipe rating sets a
minimum closure time. Computed from the Joukowsky and arithmetic-surge
relations.}
\label{fig:vol03ch11:waterhammer-surge}
\end{figure}


\paragraph{Step 5: the surge tank as the alternative mitigation.} A
slow valve is not always acceptable; an emergency trip may need to
stop the flow faster than five seconds. The standard alternative is a
surge tank near the upper end of the penstock, a vertical riser open
at the top that gives the decelerating column a free surface to push
into. When the valve shuts, the column upstream of the tank carries
on flowing into the tank, its level rising until the reservoir head
arrests it; the surge then oscillates between tank and reservoir at a
period of minutes, slow enough that the pressure rise is the modest
hydrostatic level change rather than the acoustic Joukowsky spike. The
peak level rise above the reservoir, for a frictionless surge,
follows from energy conservation of the oscillating column,
\[
z_\text{max} = u_0 \sqrt{\frac{L A}{g A_t}},
\]
where $A_t$ is the surge-tank cross-section. Sizing the tank at
$A_t = 7\,\si{\meter}^2$,
\[
z_\text{max} = 3.0\sqrt{\frac{600\times 1.131}{9.81\times 7}}
= 3.0\sqrt{\frac{678.6}{68.7}} = 3.0\times 3.14
= 9.4\,\si{\meter},
\]
a pressure rise of only $\rho g z_\text{max} = 1000\times 9.81\times
9.4 = 0.092\,\text{MPa}$ on the upstream side. The penstock below the
tank still feels the Joukowsky surge of the short reach between tank
and valve, but that reach is a small fraction of $L$, so its
round-trip time is short and even a moderately fast valve stays inside
the rating. The surge tank decouples the long upstream column from the
fast valve action; it is why almost every long high-head penstock has
one.

\paragraph{Verdict.} The bare penstock with a fast valve fails:
$3.9\,\text{MPa}$ against a $1.6\,\text{MPa}$ rating. Two mitigations
each bring it inside the rating: a governed closure of at least about
five seconds, or a surge tank that confines the fast-valve surge to a
short lower reach while limiting the upstream rise to under a tenth of
a megapascal. The decision between them is operational, the surge tank
costs a structure but permits fast emergency trips; the slow closure
costs nothing but constrains the protection scheme. The inviscid
relations of this chapter, Bernoulli for the steady head, the
Joukowsky relation for the surge, the elastic wave speed for the
timing, plus the elementary energy balance of the oscillating column,
together carry the design from a layout to a defensible verdict. None
of the analysis needed viscosity; the losses that viscosity adds
would only damp the surge and make the verdict conservative.

\subsection{The diagnostic principle for inviscid analysis}

The chapter's diagnostic question, applied to any flow under
inviscid analysis, is: do the five Bernoulli assumptions hold along
the streamline of interest, between the two points where the
equation is being applied?

\begin{itemize}[leftmargin=*]
  \item Is the flow steady? Check for time-varying boundary
    conditions, pulsations, and starting transients.
  \item Is the flow incompressible? Check the Mach number; if
    $M > 0.3$ at any point on the streamline, use the compressible
    forms of chapter 13.
  \item Is the flow inviscid along the streamline? Check whether
    the streamline passes through a boundary layer or a separated
    region; if so, Bernoulli does not apply across the viscous
    region.
  \item Is the integration along a streamline? Check whether the
    two endpoints are connected by a continuous streamline through
    a non-recirculating region.
  \item Is the flow free of heat addition, shaft work, and
    separation between the endpoints? A compressor, a turbine, a
    heater, or a separated wake between the endpoints invalidates
    the relation.
\end{itemize}

A streamline that passes all five checks is fair game for
Bernoulli; a streamline that fails any one is not. The discipline
is to identify which assumption is violated and to use the
appropriate corrective tool: the momentum equation for energy-loss
regions, the unsteady Bernoulli for time-varying flows, the
compressible analysis for high Mach, the boundary-layer equations
for wall regions, the energy equation with work and heat terms for
turbomachinery.

\begin{principle}[Bernoulli applies under five assumptions, all checked]
Bernoulli's equation $\tfrac{1}{2}\rho u^2 + p + \rho g z = p_0$
holds along a streamline of a steady, incompressible, inviscid,
non-heat-receiving, non-work-receiving flow. A calculation that
applies Bernoulli to a streamline that violates any of the five
assumptions produces a number that looks plausible and is wrong.
The diagnostic discipline is to check the assumptions before
writing the equation, and to use the momentum, energy, or
boundary-layer equations when one or more assumptions fails.
\end{principle}

Two named cases connect this chapter's failure mode to the wider
volume. The aeroelastic flutter of the original Tacoma Narrows
suspension bridge in 1940
\cite{acc:tacoma-narrows-1941}, treated as a structural-dynamics
failure in chapter 6 of this volume and as a bluff-body
aerodynamic case in chapter 13, illustrates that an inviscid
pressure analysis of a body in steady wind cannot predict the
self-excited aerodynamic moment that drove the deck to oscillate;
the unsteady, separating flow over the deck section breaks four of
the five Bernoulli assumptions at once. The common-mode failure of
all three pitot probes on Air France Flight 447 on 1 June 2009
\cite{acc:bea-af447}, treated in chapter 12 of this volume,
illustrates that the Bernoulli-based velocity measurement of a
Pitot tube requires the freestream stagnation port to be free of
ice; ice crystals at altitude blocked the ports of all three
nominally independent probes, producing wrong airspeed indications
that the autopilot recognised as inconsistent and the flight crew
could not resolve. The Bernoulli physics was not at fault; the
boundary condition on the measurement (an unobstructed stagnation
port) was. Sensor independence in design did not deliver sensor
independence in failure, and the case has become a canonical
lesson in the limits of measurement under shared environmental
cause.

\section{Project}\pathtag{core}

\begin{project}[Build and calibrate a Pitot tube; measure household fan flow]
The reader builds a simple Pitot-static probe from common
materials, calibrates it against a known reference, and uses it to
map the velocity field downstream of a household fan. The
deliverable is a twenty- to thirty-page engineering report.

\textbf{Track}. Physical build plus measurement plus analysis.
\textbf{Hazard class}: low.

\textbf{Construction}. The probe consists of a forward-facing tube
(stagnation port) and a side-vented tube (static port), each
connected to one side of a differential pressure sensor or a
U-tube manometer. Suitable materials: stainless-steel hypodermic
tubing of two diameters (the inner stagnation tube nested inside an
outer static tube with side holes), or two separate rigid plastic
tubes mounted parallel with appropriate end geometries. The static
port should be at least four tube diameters back from the leading
edge to be free of the stagnation-region pressure rise.

\textbf{Calibration}. The reader calibrates the probe against a
known flow speed. Options:

\begin{itemize}[leftmargin=*]
  \item A pendulum-style anemometer (a light flap on a hinge,
    calibrated against the deflection at a known wind), useful for
    speeds above $1\,\si{\meter\per\second}$.
  \item A free-fall calibration: the probe is mounted on a moving
    arm or a falling carriage of known velocity, the differential
    pressure recorded at the known speed.
  \item A reference from a fan's manufacturer specification at a
    known rotational speed (typically accurate to within
    $\pm 20\,\%$, sufficient for the calibration of an amateur
    probe).
\end{itemize}

The calibration produces a relation $u = K\sqrt{2 \Delta p/\rho}$,
where $K$ is a probe constant close to $1.0$ for well-built
probes. The reduction script \path{code/pitot_calibration.py} fits
$K$ from reference points by least squares, applies it to a radial
survey, and integrates the axisymmetric profile to a volume flow
rate; the worked survey it writes to \path{data/pitot_fan_survey.csv}
fits $K = 0.997$ and returns about $700\,\si{\meter\cubed}$ per hour,
a template the reader adapts to the measured fan data.

\textbf{Mapping}. The reader maps the velocity field downstream of
a small household fan (axial blade type, output diameter $\sim
100$ to $200\,\si{\milli\meter}$). The probe is positioned at a
grid of points downstream of the fan: at least three streamwise
distances (one fan-diameter downstream, three, ten) and at least
five radial positions per streamwise distance. The reader records
the dynamic pressure at each grid point and converts to local
velocity using the calibrated probe constant.

\textbf{Analysis}. The reader produces:

\begin{itemize}[leftmargin=*]
  \item A plot of velocity profile $u(r)$ at each streamwise
    distance. The expected behaviour: a peaked profile near the
    fan with a velocity defect along the axis (the hub region),
    transitioning to a more uniform profile downstream as the wake
    spreads.
  \item An estimate of the total volume flow rate $Q = \int u \,dA$
    by numerical integration of the velocity profile, compared to
    the fan's nameplate rating.
  \item A discussion of the discrepancies between the inviscid
    Bernoulli expectation (uniform downstream velocity, no
    spreading) and the measured profile, naming the viscous and
    rotational effects that the inviscid analysis omits.
\end{itemize}

\textbf{Reflection ($1{,}500$ to $3{,}000$ words)}. The reader
addresses:

\begin{itemize}[leftmargin=*]
  \item The probe's accuracy; the source of the largest
    measurement uncertainty (probe alignment, tube response time,
    sensor calibration, or otherwise).
  \item The departure of the measured flow from the inviscid
    Bernoulli prediction, named in physical terms (boundary-layer
    growth, wake spreading, hub recirculation).
  \item One design improvement to the probe and one design
    improvement to the calibration that would tighten the
    measurement, with a brief analysis of why each would help.
\end{itemize}

\textbf{Deliverable}. The probe (photographed and dimensioned),
the calibration data, the mapping data, the analysis, and the
reflection.

\textbf{Time}. Two to four hours of probe construction, two to
four hours of calibration, three to six hours of mapping, four to
eight hours of write-up.
\end{project}

\section{Exercises}

The exercises mix calculation, derivation, estimation, simulation,
design, diagnosis, failure analysis, and judgment.

\subsection*{Calculation}\pathtag{core}

\begin{exercise}
A horizontal pipe carries water at $1.5\,\si{\meter\per\second}$ in
a section of $100\,\si{\milli\meter}$ diameter. The pipe contracts
smoothly to $50\,\si{\milli\meter}$ diameter. Find the flow speed
and the pressure drop in the narrower section, assuming inviscid
flow.

\paragraph{Solution sketch.} Continuity sets the speed: the area
ratio is $(100/50)^2 = 4$, so $u_2 = 4 u_1 = 4\times 1.5 =
6\,\si{\meter\per\second}$. Bernoulli on a horizontal streamline
($z_1 = z_2$) gives the pressure drop $p_1 - p_2 = \tfrac12\rho(u_2^2
- u_1^2) = 0.5\times 1000\times(36 - 2.25) = 1.69\times
10^4\,\text{Pa} \approx 16.9\,\text{kPa}$. The narrow section runs
$16.9\,\text{kPa}$ below the wide section.
\end{exercise}

\begin{exercise}
A Pitot tube on an aircraft reads a stagnation-static differential
pressure of $5\,\text{kPa}$. The air density at the flight
altitude is $1.0\,\si{\kilogram\per\meter\cubed}$. Find the
airspeed.

\paragraph{Solution sketch.} The Pitot relation gives $u =
\sqrt{2(p_0 - p)/\rho} = \sqrt{2\times 5000/1.0} = \sqrt{10\,000} =
100\,\si{\meter\per\second}$. The reading is an indicated airspeed
based on the assumed density; at the lower true density of altitude
the true airspeed is higher than this indicated value, which is why
the density used must be the local one.
\end{exercise}

\begin{exercise}
A water tank has a small hole $2\,\si{\meter}$ below the free
surface. Find the speed of the water emerging from the hole and
the volume flow rate through a hole of $10\,\si{\milli\meter}$
diameter, assuming a discharge coefficient of $0.62$.

\paragraph{Solution sketch.} Torricelli gives the ideal jet speed
$u = \sqrt{2gh} = \sqrt{2\times 9.81\times 2} =
6.26\,\si{\meter\per\second}$. The hole area is $A = \pi(0.005)^2 =
7.85\times 10^{-5}\,\si{\meter}^2$. The actual discharge applies the
coefficient, $Q = C_d A u = 0.62\times 7.85\times 10^{-5}\times 6.26
= 3.05\times 10^{-4}\,\si{\meter\cubed\per\second} \approx
0.30\,\text{L/s}$. The coefficient absorbs the jet contraction (the
vena contracta) and the small viscous loss.
\end{exercise}

\begin{exercise}
A Venturi meter in a horizontal pipe has $100\,\si{\milli\meter}$
inlet diameter and $40\,\si{\milli\meter}$ throat diameter, with
$C_d = 0.97$. The measured static pressure difference is
$30\,\text{kPa}$, the fluid is water. Find the volume flow rate.

\paragraph{Solution sketch.} The areas are $A_1 = \pi(0.05)^2 =
7.85\times 10^{-3}\,\si{\meter}^2$ and $A_2 = \pi(0.02)^2 = 1.257
\times 10^{-3}\,\si{\meter}^2$, with area ratio $A_1/A_2 = 6.25$. The
inlet speed is $u_1 = \sqrt{2(p_1 - p_2)/[\rho((A_1/A_2)^2 - 1)]} =
\sqrt{2\times 30\,000/[1000\times(39.06 - 1)]} = \sqrt{60\,000/38\,060}
= 1.256\,\si{\meter\per\second}$. The ideal flow rate is $A_1 u_1 =
9.86\,\text{L/s}$; applying $C_d = 0.97$ gives $Q = 9.57\,\text{L/s}$,
matching \path{code/venturi_flow.py}.
\end{exercise}

\begin{exercise}
Air flows steadily over a flat plate at $50\,\si{\meter\per\second}$
in standard sea-level air ($\rho = 1.225\,\si{\kilogram\per\meter
\cubed}$). Find the dynamic pressure and the stagnation pressure
just above the plate, given an ambient static pressure of
$101.3\,\text{kPa}$.

\paragraph{Solution sketch.} The dynamic pressure is $q = \tfrac12
\rho U^2 = 0.5\times 1.225\times 50^2 = 1531\,\text{Pa} \approx
1.53\,\text{kPa}$. The stagnation pressure is the static plus the
dynamic, $p_0 = p + q = 101.3 + 1.53 = 102.8\,\text{kPa}$. At Mach
$50/340 \approx 0.15$ the incompressible form is accurate to under
$1\,\%$.
\end{exercise}

\begin{exercise}
A symmetric thin airfoil at $4^{\circ}$ angle of attack in a
$60\,\si{\meter\per\second}$ airstream ($\rho =
1.0\,\si{\kilogram\per\meter\cubed}$) has chord $1.5\,\si{\meter}$.
Compute the lift coefficient, the bound circulation per unit span,
and the lift per unit span.

\paragraph{Solution sketch.} Convert the angle: $\alpha = 4^\circ =
0.0698\,\text{rad}$. The lift coefficient is $C_L = 2\pi\alpha =
0.439$. The dynamic pressure is $q = \tfrac12\rho U^2 = 0.5\times 1.0
\times 60^2 = 1800\,\text{Pa}$, so the lift per unit span is $L' =
C_L\,q\,c = 0.439\times 1800\times 1.5 = 1185\,\si{\newton\per
\meter}$. The bound circulation follows from Kutta-Joukowski, $\Gamma
= L'/(\rho U) = 1185/(1.0\times 60) = 19.7\,\si{\meter\squared\per
\second}$, equal to $\pi c U\alpha = \pi\times 1.5\times 60\times
0.0698$ as a check.
\end{exercise}

\begin{exercise}
A water main of $300\,\si{\milli\meter}$ diameter carries water at
$2\,\si{\meter\per\second}$. The main rises $15\,\si{\meter}$
between two measurement points. Find the static pressure at the
upper point if the lower point has $400\,\text{kPa}$ static
pressure, assuming inviscid steady flow.

\paragraph{Solution sketch.} The pipe has constant diameter, so
continuity gives $u_1 = u_2 = 2\,\si{\meter\per\second}$ and the
velocity heads cancel. Bernoulli then reduces to the hydrostatic
balance $p_2 = p_1 - \rho g(z_2 - z_1) = 400\,000 - 1000\times 9.81
\times 15 = 400\,000 - 147\,150 = 253\,\text{kPa}$. The static
pressure drops by the elevation head $\rho g\,\Delta z$ alone; raising
the water $15\,\si{\meter}$ costs about $147\,\text{kPa}$.
\end{exercise}

\begin{exercise}
A circular cylinder of radius $50\,\si{\milli\meter}$ is in a
uniform freestream of $10\,\si{\meter\per\second}$ in water
($\rho = 1000\,\si{\kilogram\per\meter\cubed}$). With no
circulation, find the velocity and pressure coefficient at the top
of the cylinder. State why the integrated pressure gives zero
drag.

\paragraph{Solution sketch.} At the top, $\theta = 90^\circ$, the
surface speed is $u_\theta = -2U\sin\theta = -2\times 10\times 1 =
-20\,\si{\meter\per\second}$ (magnitude $20\,\si{\meter\per\second}$,
twice the freestream). The pressure coefficient is $C_p = 1 -
4\sin^2\theta = 1 - 4 = -3$, so the gauge pressure is $p - p_\infty =
C_p\,\tfrac12\rho U^2 = -3\times 0.5\times 1000\times 100 = -1.5\times
10^5\,\text{Pa} = -150\,\text{kPa}$. The integrated pressure gives
zero drag because $C_p(\theta) = 1 - 4\sin^2\theta$ is symmetric
front to back: the pressure on the upstream face is matched by an
equal pressure on the downstream face, so the streamwise components
cancel (d'Alembert's paradox). The radius value is not needed for
$C_p$.
\end{exercise}

\begin{exercise}
A sharp-crested rectangular weir of width $2\,\si{\meter}$ has a
head of $0.30\,\si{\meter}$ above the crest. Using $C_d = 0.62$,
find the discharge over the weir.

\paragraph{Solution sketch.} The rectangular-weir formula gives $Q =
C_d\,\tfrac23 b\sqrt{2g}\,H^{3/2} = 0.62\times 0.667\times 2\times
\sqrt{19.62}\times 0.30^{1.5}$. Evaluating, $\sqrt{19.62} = 4.43$ and
$0.30^{1.5} = 0.164$, so $Q = 0.62\times 0.667\times 2\times 4.43
\times 0.164 = 0.60\,\si{\meter\cubed\per\second}$. The $H^{3/2}$
scaling means a $10\,\%$ rise in head raises the discharge by about
$15\,\%$, which is what makes a weir a sensitive flow gauge.
\end{exercise}

\begin{exercise}
A garden hose ($20\,\si{\milli\meter}$ diameter) carries water at
$0.5\,\si{\meter\per\second}$. The end is partially covered so the
exit area is reduced by a factor of four. Find the exit jet speed,
the dynamic pressure at the exit, and the supply pressure required
upstream, assuming inviscid flow and atmospheric exit pressure.

\paragraph{Solution sketch.} Continuity with a fourfold area
reduction gives an exit speed $u_e = 4\times 0.5 =
2\,\si{\meter\per\second}$. The exit dynamic pressure is $q_e =
\tfrac12\rho u_e^2 = 0.5\times 1000\times 4 = 2000\,\text{Pa} =
2\,\text{kPa}$. Bernoulli from the hose interior (speed
$0.5\,\si{\meter\per\second}$) to the atmospheric exit gives the
required gauge supply pressure $p_s - p_\text{atm} = \tfrac12\rho
(u_e^2 - u_1^2) = 0.5\times 1000\times(4 - 0.25) =
1875\,\text{Pa}\approx 1.9\,\text{kPa}$ above atmospheric. The supply
pressure exceeds the exit dynamic pressure only by the small inlet
kinetic term.
\end{exercise}

\subsection*{Derivation}\pathtag{standard}

\begin{exercise}
Derive the differential continuity equation $\partial \rho/
\partial t + \nabla\cdot(\rho\vec{u}) = 0$ from the integral form,
stating the assumptions on the control volume.

\paragraph{Solution sketch.} Start from $\frac{d}{dt}\int_\mathcal{V}
\rho\,dV + \oint_{\partial\mathcal{V}}\rho\vec{u}\cdot\hat{n}\,dA =
0$. For a fixed (time-independent) control volume the time derivative
passes inside as a partial: $\int_\mathcal{V}\partial\rho/\partial
t\,dV$. The divergence theorem turns the surface integral into
$\int_\mathcal{V}\nabla\cdot(\rho\vec{u})\,dV$, which requires
$\rho\vec{u}$ to be continuously differentiable in $\mathcal{V}$.
Combining, $\int_\mathcal{V}[\partial\rho/\partial t +
\nabla\cdot(\rho\vec{u})]\,dV = 0$. Since the control volume is
arbitrary, the integrand must vanish pointwise, giving the
differential form. The assumptions are a fixed control volume and a
smooth flow field (no shocks or interfaces inside $\mathcal{V}$).
\end{exercise}

\begin{exercise}
Derive Bernoulli's equation $\tfrac{1}{2}\rho u^2 + p + \rho g z =
p_0$ from the steady incompressible Euler equation by integration
along a streamline. State each assumption as it is invoked.

\paragraph{Solution sketch.} Start from steady Euler, $(\vec{u}\cdot
\nabla)\vec{u} = -\tfrac1\rho\nabla p + \vec{g}$ (steady assumption
drops $\partial\vec{u}/\partial t$). Dot with the unit tangent
$\hat{s}$ and use $(\vec{u}\cdot\nabla)\vec{u}\cdot\hat{s} =
u\,\partial u/\partial s = \tfrac12\partial(u^2)/\partial s$. The
pressure term becomes $-\tfrac1\rho\partial p/\partial s$ and gravity
$-g\,\partial z/\partial s$. With $\rho$ constant (incompressible),
multiply by $\rho$ and integrate along $s$ from point 1 to point 2:
$\int(\tfrac\rho2\partial(u^2)/\partial s + \partial p/\partial s +
\rho g\,\partial z/\partial s)\,ds = 0$, yielding $\tfrac12\rho u^2 +
p + \rho g z = \text{const}$. The inviscid assumption entered when
Euler omitted viscous stress; the along-a-streamline integration
fixes the constant only on that streamline.
\end{exercise}

\begin{exercise}
Show that the velocity potential $\phi$ for an irrotational
incompressible flow satisfies Laplace's equation $\nabla^2 \phi =
0$.

\paragraph{Solution sketch.} Irrotationality means $\nabla\times
\vec{u} = 0$, which (in a simply connected region) lets the velocity
be written as a gradient, $\vec{u} = \nabla\phi$; this guarantees
$\nabla\times\nabla\phi = 0$ identically. Incompressibility gives
$\nabla\cdot\vec{u} = 0$. Substituting the first into the second,
$\nabla\cdot(\nabla\phi) = \nabla^2\phi = 0$. The two physical
conditions, no rotation and no volume change, combine into Laplace's
equation; its harmonic solutions are the same class that governs
electrostatic potential and steady heat conduction.
\end{exercise}

\begin{exercise}
Derive the Kutta-Joukowski theorem $L' = \rho U_\infty \Gamma$ by
integrating the pressure around a closed body in a uniform inviscid
flow with bound circulation $\Gamma$.

\paragraph{Solution sketch.} The compact route uses the complex
potential. Write the complex velocity far from the body as a Laurent
series $w(z) = dF/dz = U_\infty + \Gamma/(2\pi i z) + \cdots$. The
Blasius force theorem gives the force per unit span as $F_x - iF_y =
\tfrac{i\rho}{2}\oint (w)^2\,dz$. Substituting the series and keeping
the residue (the $1/z$ term survives the contour integral), the
integral picks up $2\pi i$ times the coefficient of $1/z$ in $w^2$,
which is $2U_\infty\,\Gamma/(2\pi i)$. The algebra yields $F_x = 0$
(no drag, d'Alembert) and $F_y = \rho U_\infty\Gamma$, the
Kutta-Joukowski lift. The real-variable route integrates the
Bernoulli pressure $p = p_0 - \tfrac12\rho|\vec{u}|^2$ over the
surface with the asymmetric speed produced by the added circulation;
only the cross term between freestream and circulation survives, again
giving $\rho U_\infty\Gamma$.
\end{exercise}

\begin{exercise}
Show that, for a converging-diverging nozzle in inviscid
incompressible flow, the pressure recovery is complete if the exit
area equals the inlet area. State why real diffusers do not
recover the full dynamic pressure.

\paragraph{Solution sketch.} Continuity gives $u_\text{exit} =
u_\text{inlet}(A_\text{inlet}/A_\text{exit})$; with equal areas the
exit speed equals the inlet speed. Bernoulli on the centreline (no
elevation change) then gives $p_\text{exit} = p_\text{inlet} +
\tfrac12\rho(u_\text{inlet}^2 - u_\text{exit}^2) = p_\text{inlet}$:
the pressure returns exactly. Real diffusers fall short because the
diverging section runs against an adverse pressure gradient, the
boundary layer thickens and separates, and the separated region
dissipates mechanical energy as turbulence; the lost energy is not
available to recover pressure, so a good diffuser recovers only
$70$ to $90\,\%$ of the dynamic pressure.
\end{exercise}

\begin{exercise}
Derive the unsteady Bernoulli equation $\partial \phi/\partial t +
\tfrac{1}{2}u^2 + p/\rho + gz = f(t)$ for irrotational
incompressible unsteady flow. Identify the term that vanishes
under the steady assumption.

\paragraph{Solution sketch.} For irrotational flow $\vec{u} =
\nabla\phi$. Substitute into the full Euler equation and use the
vector identity $(\vec{u}\cdot\nabla)\vec{u} = \nabla(\tfrac12 u^2) -
\vec{u}\times(\nabla\times\vec{u})$; the last term vanishes because
the flow is irrotational. The local-acceleration term is
$\partial\vec{u}/\partial t = \nabla(\partial\phi/\partial t)$. Euler
then reads $\nabla[\partial\phi/\partial t + \tfrac12 u^2 + p/\rho +
gz] = 0$, so the bracket is a function of time alone, $f(t)$. The
term that vanishes under the steady assumption is $\partial\phi/
\partial t$; dropping it recovers the ordinary Bernoulli equation
with $f(t)$ a constant.
\end{exercise}

\subsection*{Estimation}\pathtag{core}

\begin{exercise}
Estimate the lift produced by a typical paper airplane at
release. State assumptions about the airplane's wing area, angle
of attack, release speed, and air density. Compare to the
airplane's weight.

\paragraph{Reference estimate.} Take wing area $A \approx
0.03\,\si{\meter}^2$, release speed $U \approx 6\,\si{\meter\per
\second}$, angle of attack $\alpha \approx 5^\circ$, air density
$1.2\,\si{\kilogram\per\meter\cubed}$, and (for a low-aspect-ratio
flat plate, below the thin-airfoil slope) $C_L \approx 0.3$. Then $L =
C_L\,\tfrac12\rho U^2 A = 0.3\times 0.5\times 1.2\times 36\times 0.03
\approx 0.19\,\si{\newton}$. A sheet of paper massing about
$5\,\si{\gram}$ weighs $0.05\,\si{\newton}$, so the lift is a few
times the weight, consistent with a plane that glides rather than
drops. The estimate is good to a factor of two or three; the largest
uncertainty is $C_L$ at low aspect ratio.
\end{exercise}

\begin{exercise}
Estimate the maximum flow speed downstream of a household sink tap
when fully opened. State assumptions about the supply pressure,
the tap geometry, and the length of any vertical drop from supply
to spout.

\paragraph{Reference estimate.} Household mains pressure is about
$300\,\text{kPa}$ gauge ($3\,\text{bar}$). Treating the tap as an
opening to atmosphere and applying Bernoulli from the pressurised
supply to the open spout, the ideal jet speed is $u = \sqrt{2\Delta
p/\rho} = \sqrt{2\times 300\,000/1000} = \sqrt{600} \approx
24\,\si{\meter\per\second}$. Real taps run far slower because the
valve, aerator, and pipe friction drop most of the head; a measured
spout speed of $2$ to $4\,\si{\meter\per\second}$ is typical, so the
inviscid Bernoulli value is an upper bound roughly $6$ to $10$ times
the real speed. The gap is the chapter's standing warning: Bernoulli
sets the ceiling, not the operating point, when losses dominate.
\end{exercise}

\begin{exercise}
Estimate the volume flow rate through a leaking garden hose with a
pinhole at the supply (mains) end. State assumptions about the
supply pressure and the pinhole diameter; bracket the answer to
within a factor of three.

\paragraph{Reference estimate.} Take supply pressure $300\,\text{kPa}$
gauge and a pinhole of diameter $1\,\si{\milli\meter}$, area $A =
\pi(0.0005)^2 = 7.85\times 10^{-7}\,\si{\meter}^2$. The Torricelli
speed is $u = \sqrt{2\Delta p/\rho} = \sqrt{600} \approx
24\,\si{\meter\per\second}$. With a discharge coefficient $C_d
\approx 0.6$ for a sharp hole, $Q = C_d A u = 0.6\times 7.85\times
10^{-7}\times 24 \approx 1.1\times 10^{-5}\,\si{\meter\cubed\per
\second} \approx 0.7\,\text{L/min}$. Bracketing the pinhole between
$0.5$ and $2\,\si{\milli\meter}$ (a factor of four in area) spans
roughly $0.2$ to $3\,\text{L/min}$, inside the requested factor of
three about a central estimate near $0.7\,\text{L/min}$.
\end{exercise}

\begin{exercise}
Estimate the wind speed that would cause a sheet of office paper
to lift off a flat desk. State the paper's mass per unit area and
the relevant pressure-coefficient differential between upper and
lower surfaces.

\paragraph{Reference estimate.} Office paper masses about
$80\,\si{\gram\per\meter\squared}$, so its weight per unit area is
$0.08\times 9.81 \approx 0.78\,\si{\pascal}$. Lift requires a
pressure differential of this order across the sheet. A wind blowing
over the top creates a suction; take a pressure-coefficient
differential $\Delta C_p \approx 0.5$. Then $\tfrac12\rho U^2\,\Delta
C_p = 0.78$ gives $U^2 = 2\times 0.78/(1.2\times 0.5) = 2.6$, so $U
\approx 1.6\,\si{\meter\per\second}$. A gentle draught of a couple of
metres per second lifts a desk sheet, matching everyday experience;
the dominant uncertainty is $\Delta C_p$, which depends on how the
sheet sits and curls.
\end{exercise}

\begin{exercise}
Estimate the time required to drain a swimming pool of $25\,\si{
\meter} \times 10\,\si{\meter}$ surface area and $2\,\si{\meter}$
depth through a $100\,\si{\milli\meter}$ outlet near the bottom.
State assumptions about the discharge coefficient and the
behaviour as the surface drops to the outlet.

\paragraph{Reference estimate.} The closed-form drain time for a
constant-area surface is $t = (A_s/C_d A_o)\sqrt{2h_0/g}$. With $A_s =
250\,\si{\meter}^2$, $A_o = \pi(0.05)^2 = 7.85\times
10^{-3}\,\si{\meter}^2$, $C_d = 0.62$, and $h_0 = 2\,\si{\meter}$:
$\sqrt{2\times 2/9.81} = 0.639\,\si{\second}^{1/2}\cdot\si{\meter}^{1/2}$
factor, and $t = (250/(0.62\times 7.85\times 10^{-3}))\times 0.639
\approx 51\,400\times 0.639 \approx 3.3\times 10^4\,\si{\second}
\approx 9\,\text{h}$. The $\sqrt{h}$ dependence means the last shallow
layer drains slowly; once the surface nears the outlet the simple
Torricelli model breaks down (vortexing and air entrainment), so the
estimate is reliable for the bulk of the drain and optimistic for the
final centimetres. The script \path{code/tank_drain.py} confirms the
nine-hour figure.
\end{exercise}

\subsection*{Simulation}\pathtag{mastery}

\begin{exercise}
Implement a numerical solver for the two-dimensional Laplace
equation $\nabla^2 \phi = 0$ on a rectangular domain with mixed
Dirichlet and Neumann boundary conditions, using finite differences.
Verify on the flow around a circular cylinder by comparing the
computed surface pressure coefficient to the analytical $C_p =
1 - 4\sin^2\theta$.

\paragraph{Solution sketch.} Discretise the domain on a grid and
apply the five-point Laplacian; Gauss-Seidel or successive
over-relaxation iterates $\phi_{i,j} = \tfrac14(\phi_{i+1,j} +
\phi_{i-1,j} + \phi_{i,j+1} + \phi_{i,j-1})$ to convergence. Impose
$\phi = U_\infty x$ on the far field (Dirichlet) and zero normal
derivative on the cylinder surface (Neumann, no penetration). Recover
surface velocity by differencing $\phi$ tangentially, then $C_p = 1 -
(u_s/U_\infty)^2$. The computed distribution should reproduce the
peak $C_p = 1$ at the stagnation points and $C_p = -3$ at the top and
bottom; \path{code/cylinder_potential.py} gives the analytical target
to compare against. Convergence and grid-refinement checks are part
of the deliverable.
\end{exercise}

\begin{exercise}
Implement a point-vortex panel method for thin airfoils: discretise
the chord into panels, place a point vortex at the quarter-chord
of each panel, enforce the flow-tangency condition at the
three-quarter-chord of each panel, and solve the resulting linear
system. Verify on a flat plate by comparing the computed lift
slope to $C_L = 2\pi\alpha$.

\paragraph{Solution sketch.} Lay $N$ panels along the chord; place a
point vortex of unknown strength $\gamma_j$ at each panel's
quarter-chord and a collocation point at each three-quarter-chord.
The flow-tangency condition at each collocation point gives one
equation: the sum of induced normal velocities from all vortices plus
the freestream normal component $U_\infty\sin\alpha$ must vanish. This
is an $N\times N$ linear system $\mathbf{A}\boldsymbol{\gamma} =
\mathbf{b}$; solve it, sum the circulations $\Gamma = \sum\gamma_j$,
and form $C_L = \Gamma/(\tfrac12 U_\infty c)$. For a flat plate the
result converges to $C_L = 2\pi\alpha$ as $N$ grows; even $N = 1$ with
the quarter/three-quarter placement gives the exact slope, which is
the classical reason for that placement (the lumped-vortex theorem).
\end{exercise}

\begin{exercise}
Simulate the streaklines and pathlines of an unsteady flow field
$u_x = U_0 \cos(\omega t)$, $u_y = 0$, by integrating particle
trajectories. Show that the streamlines coincide with the pathlines
only when the flow is steady.

\paragraph{Solution sketch.} For the given field the pathline of a
particle from $(x_0,0)$ integrates to $x(t) = x_0 + (U_0/\omega)
\sin(\omega t)$, $y(t) = 0$: the particle oscillates back and forth
on a horizontal line. An instantaneous streamline at time $t_0$
follows $dx/u_x = dy/u_y$ with $u_y = 0$, so streamlines are also
horizontal lines, but their direction reverses with the sign of
$\cos(\omega t_0)$. A streakline of particles released continuously
from one point spreads along the same line with a phase-dependent
extent. With $u_y \ne 0$ added (as in the figure of this chapter) the
three curves separate visibly. The verification is to integrate
particle trajectories numerically and overlay them on the snapshot
streamlines: they coincide only when $\omega = 0$ (steady flow), and
the code should demonstrate the divergence for $\omega \ne 0$.
\end{exercise}

\subsection*{Design}\pathtag{standard}

\begin{exercise}
Design a Venturi flow meter to measure flow rates between $1$ and
$10\,\text{L/s}$ in a $50\,\si{\milli\meter}$ water pipe, with a
differential pressure transducer capable of resolving pressures
between $100\,\text{Pa}$ and $50\,\text{kPa}$. Specify the throat
diameter and check that the resulting differential pressures fall
within the transducer's range.

\paragraph{Model-answer sketch.} The pipe area is $A_1 = \pi(0.025)^2
= 1.96\times 10^{-3}\,\si{\meter}^2$. The differential pressure scales
as $\Delta p \approx \tfrac12\rho(Q/A_2)^2$ for a strong contraction.
Try a throat of $25\,\si{\milli\meter}$ ($A_2 = 4.91\times
10^{-4}\,\si{\meter}^2$, area ratio $4$): at $Q = 1\,\text{L/s}$, $u_2
= Q/A_2 = 2.04\,\si{\meter\per\second}$ and $\Delta p \approx
\tfrac12\rho(u_2^2 - u_1^2) \approx 2.0\,\text{kPa}$; at $Q =
10\,\text{L/s}$, $u_2 = 20.4\,\si{\meter\per\second}$ and $\Delta p
\approx 200\,\text{kPa}$, which exceeds the $50\,\text{kPa}$ ceiling.
Open the throat to $32\,\si{\milli\meter}$ (area ratio $\approx 2.4$):
the differential drops by roughly the square of the ratio change,
landing the $10\,\text{L/s}$ point near $45\,\text{kPa}$ and the
$1\,\text{L/s}$ point near $450\,\text{Pa}$, both inside the
transducer band. The design is iterative: pick the throat that keeps
$\Delta p$ inside the range across the full flow span, then verify
with \path{code/venturi_flow.py}.
\end{exercise}

\begin{exercise}
Design a sharp-crested rectangular weir to measure flows from a
small stream that varies between $0.05$ and $1\,\text{m}^3\text{
/s}$. Specify the crest width and the expected head range; check
that the discharge coefficient remains within typical limits over
the range.

\paragraph{Model-answer sketch.} Invert $Q = C_d\tfrac23 b\sqrt{2g}\,
H^{3/2}$ for $H$. Choose a crest width $b = 1\,\si{\meter}$. At $Q =
0.05\,\si{\meter\cubed\per\second}$: $H^{3/2} = Q/(C_d\tfrac23 b
\sqrt{2g}) = 0.05/(0.62\times 0.667\times 1\times 4.43) = 0.0273$, so
$H = 0.0273^{2/3} \approx 0.091\,\si{\meter}$. At $Q =
1\,\si{\meter\cubed\per\second}$: $H^{3/2} = 0.546$, $H \approx
0.66\,\si{\meter}$. The head ranges from about $90\,\si{\milli\meter}$
to $0.66\,\si{\meter}$, a comfortable measurement span. Check that the
low end keeps $H$ above roughly $30\,\si{\milli\meter}$ (below which
surface tension makes $C_d$ unreliable); $90\,\si{\milli\meter}$ is
safe. A wider crest lowers the heads and may push the low flow below
that limit, so $b = 1\,\si{\meter}$ is a reasonable choice.
\end{exercise}

\begin{exercise}
Design a household-scale wind-turbine blade of $2\,\si{\meter}$
diameter for operation in $5\,\si{\meter\per\second}$ winds at sea
level. Specify the blade's chord, twist, and angle of attack
distribution along the span, using thin-airfoil theory and the
inviscid Kutta-Joukowski relation. State which inviscid
assumptions break down and where, and estimate the resulting
viscous correction.

\paragraph{Model-answer sketch.} Pick a tip-speed ratio $\lambda
\approx 6$, so the tip speed is $\lambda V = 30\,\si{\meter\per
\second}$ and the local blade speed at radius $r$ is $\Omega r =
\lambda V r/R$. The local relative wind angle is $\phi = \arctan(V/
\Omega r)$, large near the root and small near the tip, which sets the
twist. Choose a working lift coefficient $C_L \approx 0.8$ at angle of
attack $\alpha \approx 5^\circ$, so the geometric pitch is $\phi -
\alpha$. The chord from blade-element-momentum theory scales roughly
as $c(r) \propto 1/(\lambda^2 r)$, giving a wide root tapering to a
narrow tip. The inviscid Kutta-Joukowski relation sets the lift per
element; the inviscid assumptions break down in the root region
(thick sections, low local Reynolds number, three-dimensional
flow) and near stall at low wind speeds. The viscous correction is a
drag-to-lift ratio of a few percent that lowers the power coefficient
below the Betz limit; expect a realised $C_P \approx 0.35$ to $0.45$.
\end{exercise}

\begin{exercise}
Design a converging-diverging diffuser to recover pressure from a
$30\,\si{\meter\per\second}$ air jet exiting a duct, with an exit
area double the inlet area. Specify the diffuser length and the
divergence angle; estimate the achievable pressure-recovery
coefficient.

\paragraph{Model-answer sketch.} Doubling the area halves the speed,
so the ideal pressure recovery is $\Delta p = \tfrac12\rho(u_1^2 -
u_2^2) = \tfrac12\rho u_1^2(1 - \tfrac14) = \tfrac34 q_1$, where $q_1
= \tfrac12\times 1.2\times 30^2 = 540\,\text{Pa}$, giving an ideal
$\Delta p = 405\,\text{Pa}$. To avoid separation keep the total
included divergence angle at or below about $7^\circ$ (each wall
$3.5^\circ$). For a circular duct, doubling the area multiplies the
diameter by $\sqrt2 = 1.41$; the length needed for a $3.5^\circ$
half-angle is $L = (d_2 - d_1)/(2\tan 3.5^\circ) \approx 0.41 d_1/
0.122 \approx 3.4\,d_1$. The real pressure-recovery coefficient
$C_p = \Delta p/q_1$ reaches about $0.6$ to $0.7$ of the ideal $0.75$
for a well-designed diffuser, so expect $C_p \approx 0.5$.
\end{exercise}

\subsection*{Diagnosis and reverse engineering}\pathtag{core}

\begin{exercise}
A homeowner reports that water pressure at the upper-floor shower
drops sharply when the kitchen tap is opened. Identify the
inviscid-fluid-mechanics calculation that quantifies the drop and
the design discipline that would have anticipated it. State the
diagnostic measurement.

\paragraph{Solution sketch.} Opening the kitchen tap adds flow to the
shared supply pipe; continuity raises the velocity in the common run,
and Bernoulli plus pipe-friction losses convert that higher velocity
into a static-pressure drop felt at the upstream branch to the shower.
The calculation is a head balance on the shared pipe: the head loss
scales with the square of the total flow, so the additional draw drops
the pressure available at the upper floor. The design discipline is
hydraulic network sizing, oversizing the trunk and branching close to
the source so that fixtures do not load a common small-diameter run.
The diagnostic measurement is a pressure gauge at the shower, read
with the kitchen tap shut and again with it open; the difference
localises the drop to the shared run.
\end{exercise}

\begin{exercise}
A Pitot tube on a small aircraft reads an airspeed $20\,\%$ lower
than the GPS-derived groundspeed. Possible causes: probe misalignment,
static-port obstruction, calibration drift, ice. Identify which of
these the diagnostic discipline would investigate first, and the
test that would distinguish among them.

\paragraph{Solution sketch.} Investigate the static-port obstruction
first: a partially blocked static port is the single fault that biases
the dynamic-pressure reading directly and is common and cheap to
check, whereas calibration drift is slow and probe misalignment
usually produces a smaller error. A $20\,\%$ low airspeed corresponds
to a dynamic-pressure error of about $36\,\%$ (since $u\propto
\sqrt{\Delta p}$), consistent with a static port reading too high.
The distinguishing test is to compare indicated airspeed against GPS
groundspeed on reciprocal headings in still air (averaging out wind),
and to inspect and clear the static port; if clearing the port
restores agreement, the cause is confirmed. Ice is distinguished by
its dependence on conditions (it appears in cloud or at low
temperature and clears with heat).
\end{exercise}

\begin{exercise}
A small wind turbine in steady $7\,\si{\meter\per\second}$ wind is
producing only $40\,\%$ of its nameplate power. Identify three
inviscid-flow-related explanations and the inspection that would
distinguish among them.

\paragraph{Solution sketch.} Three inviscid-flow explanations: (1) the
blades operate off their design tip-speed ratio, so the local angle of
attack is wrong and the bound circulation, hence lift, is low; (2) the
blade pitch or twist is misset, again moving the angle of attack off
the design point; (3) blade fouling or leading-edge damage alters the
effective camber and the Kutta condition, reducing circulation. The
distinguishing inspections: measure the rotor speed to compute the
actual tip-speed ratio and compare to design (explanation 1); check
the pitch setting against the design schedule (explanation 2); inspect
the leading edges and surfaces for fouling, ice, or erosion
(explanation 3). Power scales with the cube of the effective wind and
strongly with $C_P(\lambda)$, so a modest tip-speed-ratio error
explains a large power shortfall.
\end{exercise}

\begin{exercise}
A dam spillway has shown signs of cavitation damage on its concrete
face downstream of the crest. Identify the inviscid-pressure
calculation that would have predicted the local pressure drop and
the corresponding cavitation risk.

\paragraph{Solution sketch.} The calculation is the Bernoulli static
pressure along the spillway face: $u = \sqrt{2g(z_1 - z_2)}$ gives the
local speed and $p = p_\text{atm} - \tfrac12\rho u^2$ (plus any
$C_p$ deficit over surface irregularities) gives the local static
pressure. Where this drops below the vapour pressure $p_v \approx
2.3\,\text{kPa}$ absolute, the cavitation number $\sigma = (p -
p_v)/\tfrac12\rho u^2$ falls below inception and bubbles form, then
collapse downstream and pit the concrete (see
Figure~\ref{fig:vol03ch11:cavitation} and the
$30\,\si{\meter}$-spillway estimation). The discipline is to compute
$\sigma$ at every surface irregularity for the design discharge and,
where it is marginal, to add aeration; the diagnostic confirmation is
that damage clusters just downstream of joints and offsets where the
local pressure dips.
\end{exercise}

\subsection*{Failure analysis}\pathtag{core}

\begin{exercise}
The Air France Flight 447 accident report \cite{acc:bea-af447}
identifies pitot-icing-induced unreliable airspeed as a precipitating
condition. Argue, in $200$ to $300$ words, what the inviscid
Bernoulli physics of the Pitot tube does and does not guarantee
about the airspeed measurement, and what design discipline would
reduce the probability of a common-mode failure across three
nominally independent probes.

\paragraph{Model-answer sketch.} The Bernoulli physics guarantees
that, given an unobstructed stagnation port and an accurate static
pressure, the differential $p_0 - p$ equals the dynamic pressure
$\tfrac12\rho u^2$, so the airspeed is recoverable. It guarantees
nothing about the boundary condition on the measurement: it assumes
the stagnation port samples the true total pressure. Ice crystals
blocking the ports break that assumption, and the equation then reports
a wrong speed that looks valid. A strong answer notes that the three
probes were independent in wiring and electronics but shared one
environment, so an environmental cause (high-altitude ice) defeats the
redundancy: independence in design is not independence in failure
under a common cause. The design discipline is diversity, not just
redundancy: probes of differing heating power, geometry, or operating
principle, plus a model-based airspeed estimate from other sensors to
cross-check, so that no single environmental condition disables all
channels at once.
\end{exercise}

\begin{exercise}
A student applies Bernoulli's equation to a streamline that enters
a turbomachine inlet and exits the turbomachine outlet, getting an
inconsistent answer. Identify which of the five Bernoulli
assumptions is violated and the correct analytical tool.

\paragraph{Solution sketch.} The violated assumption is no shaft work
between the endpoints: a turbomachine adds or extracts mechanical
energy across the control volume, so the Bernoulli sum is not constant
along a streamline through the machine. The correct tool is the steady
energy equation with a work term, $\tfrac{p_1}\rho + \tfrac12 u_1^2 +
g z_1 + w_\text{shaft} = \tfrac{p_2}\rho + \tfrac12 u_2^2 + g z_2 +
\text{losses}$, where $w_\text{shaft}$ is the specific work of the
machine. Bernoulli is the special case with $w_\text{shaft} = 0$ and
zero losses; restoring the work term resolves the inconsistency.
\end{exercise}

\begin{exercise}
A garden-hose nozzle produces a jet whose computed Bernoulli speed
disagrees with the measured speed by a factor of two. Identify the
two most probable causes (one in the assumption set, one in the
measurement) and the diagnostic that would distinguish them.

\paragraph{Solution sketch.} Assumption-set cause: the nozzle is not
inviscid and not loss-free; pipe friction and the nozzle contraction
drop much of the supply head, so the real exit speed is well below the
Bernoulli value computed from the nominal supply pressure (the same
ceiling-not-operating-point gap as the sink-tap estimate). Measurement
cause: the speed was inferred from an indirect quantity (throw
distance, or a pressure read at the wrong point) rather than measured
directly. The diagnostic is a direct measurement: collect a known
volume over a timed interval to get $Q$, divide by the measured exit
area for the true mean speed, and measure the static pressure at the
nozzle inlet rather than at the supply; comparing the head actually
available at the nozzle to the nominal supply pressure shows whether
the loss (assumption) or the inference (measurement) is the larger
error.
\end{exercise}

\subsection*{Judgment}\pathtag{standard}

\begin{exercise}
A senior engineer argues that Bernoulli's equation is sufficient
for any subsonic-flow design problem because viscous corrections
are typically a few percent. Write a one-paragraph reply that
agrees where agreement is honest and disagrees where the
viscous-correction estimate fails.

\paragraph{Model-answer sketch.} Agree that for attached external
flow at moderate angle of attack, and for pressure recovery in
gently varying ducts, viscosity costs only a few percent and Bernoulli
gives the right leading-order answer; that is the working basis of
subsonic aerodynamics and flow metering. Disagree where the
viscous-correction estimate is not small: separated flow (the entire
drag of a bluff body, stall, diffusers past the divergence limit),
flows through boundary layers, hydraulic jumps, and any
energy-dissipating transition. In those cases the correction is not a
few percent but order one or larger, and Bernoulli is qualitatively
wrong; a senior engineer should treat the few-percent rule as valid
only after confirming the flow is attached and loss-free between the
points of interest.
\end{exercise}

\begin{exercise}
A junior engineer is taught the equal-transit-time explanation of
airfoil lift in an introductory physics course. Write a
one-paragraph response that identifies the explanation's errors
and states the correct physics in terms of bound circulation and
the Kutta condition.

\paragraph{Model-answer sketch.} The equal-transit-time story assumes
fluid particles split at the leading edge must rejoin at the trailing
edge; there is no such constraint, and on a real airfoil the upper
flow arrives well ahead of the lower flow. It also predicts no lift
for a flat plate or a symmetric airfoil, which plainly do lift at
angle of attack, and predicts negative lift for some cambered shapes
that lift. The correct physics: a real airfoil at angle of attack
develops a bound circulation $\Gamma$ fixed by the Kutta condition
(the flow leaves the sharp trailing edge smoothly). The circulation
adds to the freestream, producing a faster upper flow and a slower
lower flow; Bernoulli converts the speed difference into a pressure
difference, and the integrated difference is the lift $L' = \rho
U_\infty\Gamma$. The path length plays no role; the trailing-edge
condition and the resulting circulation do.
\end{exercise}

\begin{exercise}
A designer proposes to compute the drag of a road vehicle by
integrating the inviscid Bernoulli pressure distribution over the
vehicle's surface. Write a one-paragraph response covering why
this approach gives zero drag (d'Alembert's paradox), what
modification recovers a sensible drag estimate, and what physical
mechanism the inviscid analysis omits.

\paragraph{Model-answer sketch.} Integrating the inviscid Bernoulli
pressure over a closed body gives zero net streamwise force because
the surface pressure distribution is front-to-back symmetric: the
high pressure on the front is exactly recovered as high pressure on
the rear (d'Alembert's paradox, shown for the cylinder by $C_p = 1 -
4\sin^2\theta$). The modification that recovers a sensible estimate is
to include the wake: a real bluff body separates, the rear pressure
never recovers, and the front-to-rear pressure imbalance is the
pressure (form) drag, conventionally captured by a measured drag
coefficient $C_D$ so that $D = C_D\,\tfrac12\rho U^2 A$. The omitted
mechanism is viscosity, which creates the boundary layer, causes
separation, and (through skin friction) adds a second drag component;
the inviscid model has neither.
\end{exercise}

\begin{exercise}
Identify a fluid-flow situation in your immediate environment in
which Bernoulli's equation gives a useful first estimate, and one
in which it gives a misleading answer. State, in $200$ to $300$
words, the streamline you would draw in each case and the
assumption that fails in the misleading case.

\paragraph{Solution sketch.} A strong answer pairs a clean case with
a trap. Useful case: a partially covered garden hose or a tap aerator,
where a streamline from the supply to the open jet gives a sensible
exit-speed estimate because the flow is steady, nearly incompressible,
and short enough that losses are modest. Misleading case: a streamline
drawn through a hand-dryer or vacuum-cleaner nozzle and into the
turbulent room jet, or across a window in a gust, where the flow is
unsteady and separated; applying Bernoulli across the separated wake
violates the inviscid (and often the steady) assumption and gives a
wrong pressure. The discipline shown is to state the streamline
explicitly and check the five assumptions on it before trusting the
number.
\end{exercise}
